% !TeX root = ../../Book1.tex
\chapter{Poincaré 不等式与 Sobolev 嵌入}
\label{b1:ch:24}
\BookChapterNumber{b1:ch:24}
\begin{chapterabstract}
本章研究弱梯度能够控制哪些函数性质。减去均值或施加零边界条件后，Poincaré 不等式控制函数振幅；当导数可积指数低于维数时，Sobolev 不等式提高函数的可积指数；超过维数时，Morrey 估计给出具有确定 Hölder 指数的连续代表。高阶结论通过迭代与局部化建立，临界情形则转向有限指数常数的增长和指数可积性。不同结论对区域几何、边界条件与指数端点各有要求。
\end{chapterabstract}

\begin{learninggoals}
\begin{enumerate}
\item\label{b1:goal:24-poincare}证明球上的均值估计、任意有界域的零边界估计及连通 Lipschitz 区域的均值估计，明确常数核和几何常数的作用。
\item\label{b1:goal:24-sobolev}从逐坐标积分与任意维乘积估计出发证明 Gagliardo--Nirenberg--Sobolev 不等式，再用幂次、稠密性和延拓处理一般指数与区域。
\item\label{b1:goal:24-morrey}通过逐尺度均值比较在每个点构造连续代表，推导局部、全空间及闭区域上的 Morrey 估计，并辨认指数的适用边界。
\item\label{b1:goal:24-higher}把一阶结论组织成高阶和分数阶嵌入，区分经典连续导数、弱导数与高阶边界延拓的额外要求。
\item\label{b1:goal:24-critical}在选读部分追踪临界 \(L^q\) 常数随 \(q\) 的增长，证明非最优常数的 Trudinger--Moser 不等式，并用具体函数辨认临界有界性的失败。
\end{enumerate}
\end{learninggoals}

本章以\chapref{b1:ch:21}的弱导数运算和代表理论为基础，反复使用\chapref{b1:ch:22}的内部逼近、零边界闭包及一阶延拓。\chapref{b1:ch:23}的分数阶空间用于高阶一节末尾，迹则帮助解释零边界条件。这里先直接证明均值 Poincaré 不等式，不把下一章的紧嵌入提前作为工具。

读到每个估计时，宜先检查右端是否包含函数的零阶范数，以及常数依赖于什么。全空间的紧支集路线能给出只含梯度的估计；一般有界区域若不减均值、不施加零边界条件，就必须保留常数函数的贡献。次临界可积性和超临界连续性虽都来自一阶导数，却不能通过简单替换指数相互推导。

中文伴读可参阅王明新的《索伯列夫空间》\cite[第2.6--2.11节]{Wang2013SobolevChinese}。Nirenberg 的《On elliptic partial differential equations》\cite[Lecture II]{Nirenberg1959Elliptic}提供 Sobolev 不等式的原始证明路线；本章展开其中任意维 Hölder 归纳和幂次代换。Morrey 部分参照 Hajłasz 的《Sobolev spaces on metric-measure spaces》\cite{Hajlasz2003SobolevMetric}中球均值的组织方式，但这里的证明始终在欧氏空间中完成。

% !TeX root = ../../Book1.tex
\section{Poincaré 型不等式}
\label{b1:sec:2401}

梯度不能区分两个相差常数的函数。因此，在一般函数空间中用梯度控制函数本身，必须先消除常数这一自由度。减去平均值是一种办法，要求零边界条件是另一种办法；它们需要的几何假设并不相同。本节先在球上建立尺度明确的估计，再说明连通性和边界条件分别在哪里起作用。

\subsection{平均值版本}
\label{b1:subsec:240101}

对有限正测度的集合 \(A\)，记
\[
 u_A=\frac1{|A|}\int_Au(x)\,\mathrm dx.
\]
本章 \(|A|\) 表示 Lebesgue 测度，\(|\nabla u|\) 表示弱梯度的欧氏长度。后者的 \(L^p\) 范数与本书逐偏导数求和的范数只相差维数常数。

\ProseParagraphBreak

\term[poincare-budengshi]{Poincaré 不等式}[Poincaré inequality]的球上版本如下。

\begin{theorem}[Poincaré]\label{b1:thm:ball-poincare-lp}
设 \(B=B(x_0,r)\subset\mathbb R^d\)、\(1\le p\le\infty\)。对 \(u\in W^{1,p}(B)\)，有
\begin{equation}\label{b1:eq:ball-poincare-lp}
 \|u-u_B\|_{L^p(B)}
 \le C_d r\|\nabla u\|_{L^p(B)}.
\end{equation}
常数可以只依赖于维数，不依赖于球的位置、半径或指数 \(p\)。
\end{theorem}
\begin{proof}
先取 \(u\in C^\infty(B)\cap W^{1,p}(B)\)。球的凸性保证两点间的线段留在球内，故
\[
 |u(x)-u(y)|
 \le\int_0^1|\nabla u(x+t(y-x))|\cdot|x-y|\,\mathrm dt.
\]
对 \(y\) 平均并用积分三角不等式，得到
\[
 |u(x)-u_B|
 \le\frac1{|B|}\int_0^1\int_B
   |\nabla u(x+t(y-x))|\cdot|x-y|\,\mathrm dy\,\mathrm dt.
\]
对固定 \(x,t\) 作代换 \(z=x+t(y-x)\)。Jacobi 因子为 \(t^d\)，而 \(|x-y|=|x-z|/t\)。若 \(z\in x+t(B-x)\)，则 \(|x-z|<2rt\)，因此
\[
 \begin{aligned}
 |u(x)-u_B|
 &\le\frac1{|B|}\int_B|\nabla u(z)|\cdot|x-z|
       \int_{|x-z|/(2r)}^1t^{-d-1}\,\mathrm dt\,\mathrm dz\\
 &\le C_d\int_B\frac{|\nabla u(z)|}{|x-z|^{d-1}}\,\mathrm dz.
 \end{aligned}
\]
所有被积函数非负，故换元后的次序交换由 Tonelli 定理合法；对角线上的值不影响积分。这里
\[
 \int_{a}^1t^{-d-1}\,\mathrm dt\le\frac{a^{-d}}d,
 \quad |B|=\omega_dr^d,
\]
所以半径在系数中相消。

在域外将 \(|\nabla u|\) 补零，右端受卷积核
\[
 K_r(h)=|h|^{1-d}\mathbf1_{\{|h|<2r\}}
\]
控制。这个核可积，并满足 \(\|K_r\|_1=C_dr\)。卷积的 \(L^1*L^p\) 估计于是给出\eqref{b1:eq:ball-poincare-lp}，其中常数不依赖 \(p\)。

若 \(p<\infty\)，由 Meyers--Serrin 定理在 \(W^{1,p}(B)\) 中逼近一般 \(u\)。均值是 \(L^p(B)\) 上的连续线性泛函，所以函数、均值和梯度都可取极限，得到一般情形。若 \(p=\infty\)，则 \(u\in W^{1,q}(B)\) 对每个有限 \(q\) 成立。应用刚证且常数与 \(q\) 无关的估计，再令 \(q\to\infty\)，利用有限测度集上 \(L^q\) 范数趋于本性上确界，得到无穷端点。
\end{proof}

这个证明同时给出一个有用的凸集版本。若 \(G\) 是非空有界凸开集，直径为 \(D\)，则同样的线段换元把系数换成 \(C_dD^d/|G|\)。截断核的 \(L^1\) 范数再贡献 \(C_dD\)，因此
\begin{equation}\label{b1:eq:convex-poincare-shape}
 \|u-u_G\|_{L^p(G)}
 \le C_d\frac{D^{d+1}}{|G|}\|\nabla u\|_{L^p(G)}.
\end{equation}
这里给出的是充分使用的粗常数，不主张它对所有凸集都最优。对长宽比固定的盒子，\(|G|\) 与 \(D^d\) 可比，便仍得到“直径乘梯度”的尺度。

\ProseParagraphBreak

平均值不是随意插入的修饰。若改成固定常数零，非零常函数会立即反驳估计；若把半径因子删去，伸缩 \(u(x)=v((x-x_0)/r)\) 又会使函数和梯度两边的尺度失配。估计既消除了梯度的常数核，也保留了空间尺度。

\subsection{零边界版本}
\label{b1:subsec:240102}

零边界条件已经排除非零常数，因此可以直接估计函数本身。这里不要求开集连通，也不要求它的边界具有图表。

\begin{theorem}[零边界 Poincaré 不等式]\label{b1:thm:zero-boundary-poincare}
设非空开集 \(\Omega\subset(a,b)\times\mathbb R^{d-1}\)，其中 \(b-a<\infty\)。则对 \(1\le p\le\infty\)、\(u\in W_0^{1,p}(\Omega)\)，有
\begin{equation}\label{b1:eq:zero-boundary-poincare}
 \|u\|_{L^p(\Omega)}
 \le(b-a)\|\partial_1u\|_{L^p(\Omega)}
 \le(b-a)\|\nabla u\|_{L^p(\Omega)}.
\end{equation}
特别地，每个有界开集都满足相应的零边界估计。
\end{theorem}
\begin{proof}
先取 \(\varphi\in C_c^\infty(\Omega)\)，并把它补零到全空间。对固定的其余坐标 \(x'\)，
\[
 \varphi(x_1,x')=\int_a^{x_1}\partial_1\varphi(t,x')\,\mathrm dt.
\]
有限 \(p\) 时，Hölder 不等式给出
\[
 |\varphi(x_1,x')|^p
 \le(b-a)^{p-1}\int_a^b|\partial_1\varphi(t,x')|^p\,\mathrm dt.
\]
先对 \(x_1\in(a,b)\) 积分，再对 \(x'\) 积分，所得即为\eqref{b1:eq:zero-boundary-poincare}的 \(p\) 次方。对于 \(p=1\)，沿线的积分三角不等式就是同一计算；对于 \(p=\infty\)，直接对沿线表示取上确界即可。

一般 \(u\in W_0^{1,p}(\Omega)\) 按定义有内部紧支集光滑逼近。将刚证估计传给这列逼近的极限，即得结论，包括按闭包定义给定的无穷端点。最后，\(|\partial_1u|\le|\nabla u|\) 给出第二个不等式。
\end{proof}

这一计算在\cref{b1:ex:22-strip-poincare}中已作为工具出现，现在将它列为后续反复使用的定理。它证明梯度半范数在 \(W_0^{1,p}\) 上是范数，并与原 Sobolev 范数等价。在有界 Lipschitz 区域和有限 \(p\) 时，可用上一章的零迹刻画把条件改写成 \(\operatorname{Tr}u=0\)；对任意开集或无穷端点，仍应保留闭包定义。

\subsection{Poincaré–Wirtinger 不等式}
\label{b1:subsec:240103}

在一般连通区域上，均值版本通常称为\term[poincare-wirtinger-budengshi]{Poincaré--Wirtinger 不等式}[Poincaré--Wirtinger inequality]。连通性阻止函数在不同分支上分别取互不相同的常数；边界图表则提供有限个可比较的局部区域。

\begin{theorem}[Poincaré--Wirtinger]\label{b1:thm:connected-domain-poincare}
设 \(\Omega\subset\mathbb R^d\) 为非空、有界、连通的 Lipschitz 区域。对 \(1\le p\le\infty\)，存在常数 \(C_{\Omega,p}\)，使
\[
 \|u-u_\Omega\|_{L^p(\Omega)}
 \le C_{\Omega,p}\|\nabla u\|_{L^p(\Omega)}
 \quad\bigl(u\in W^{1,p}(\Omega)\bigr).
\]
\end{theorem}
\begin{proof}
先构造有限局部覆盖。边界附近取\secref{b1:sec:2203}的有限内图表；它们与 \(\Omega\) 的交都是正半盒在双 Lipschitz 映射下的像。图表未覆盖的剩余部分是内部紧集，再用有限个内部球覆盖。由此得到有限个非空开集 \(G_1,\ldots,G_N\)，覆盖 \(\Omega\)。

每个 \(G_j\) 上都有
\[
 \|u-u_{G_j}\|_{L^p(G_j)}
 \le C_j\|\nabla u\|_{L^p(G_j)}.
\]
对于球，这是前面的定理。对于图表像，将函数拉回凸正半盒，使用\eqref{b1:eq:convex-poincare-shape}，再通过双 Lipschitz 换元控制函数和梯度。拉回的均值是一个常数 \(c_j\)，虽未必等于物理区域均值，但
\[
 \|u-u_{G_j}\|_{L^p(G_j)}
 \le2\|u-c_j\|_{L^p(G_j)}.
\]
这个估计来自
\(|u_{G_j}-c_j|\le |G_j|^{-1/p}\|u-c_j\|_p\)，无穷端点按 \(1/p=0\) 解释。故局部均值估计确实成立。

以 \(G_j\) 为顶点，若 \(G_i\cap G_j\ne\varnothing\) 就连一条边。这个有限图连通；否则两组顶点的并会把连通的 \(\Omega\) 分成两个不交非空开集。每个非空交集也是开集，测度严格为正。记 \(m_j=u_{G_j}\)，对相邻顶点及 \(A_{ij}=G_i\cap G_j\)，有
\[
 |m_i-m_j|
 \le |A_{ij}|^{-1/p}
       \bigl(\|u-m_i\|_{L^p(G_i)}
              +\|u-m_j\|_{L^p(G_j)}\bigr).
\]
这是把常数 \(m_i-m_j=(m_i-u)+(u-m_j)\) 限制到交集后取范数。

在有限连通图中固定从顶点1到各顶点的路径，沿路径相加得到
\[
 |m_j-m_1|\le C_{\Omega,p}\|\nabla u\|_{L^p(\Omega)}.
\]
常数包含局部图表界、交集测度及有限条路径。再在每个 \(G_j\) 上写
\[
 \|u-m_1\|_{L^p(G_j)}
 \le\|u-m_j\|_{L^p(G_j)}+|G_j|^{1/p}|m_j-m_1|.
\]
有限覆盖允许将这些局部估计合为
\(\|u-m_1\|_{L^p(\Omega)}\le C_{\Omega,p}\|\nabla u\|_p\)。
最后用
\(\|u-u_\Omega\|_p\le2\|u-m_1\|_p\)，完成证明。上述有限和及交集估计也适用于 \(p=\infty\)。
\end{proof}

这个证明没有用到下一章的紧嵌入定理。相反，它具体显示了常数如何受几何影响：图表越来越窄或相邻块的交集越来越小，所给的统一常数都可能恶化。仅知道区域有界，并不足以对一族变化的区域取得相同常数。

\ProseParagraphBreak

若 \(\Omega\) 不连通，取函数在两个分量上分别为不同常数，则弱梯度为零，减去全域均值后却不为零。因此通常只能逐连通分量减去各自均值，或另加能连接这些自由常数的约束。零边界版本没有这一障碍，因为每个分量上的非零常数都已被闭包条件排除。

\subsection{局部形式}
\label{b1:subsec:240104}

将球上不等式除以 \(|B|^{1/p}\)，得到尺度归一化的形式
\begin{equation}\label{b1:eq:poincare-average-form}
 \left(\frac1{|B|}\int_B|u-u_B|^p\,\mathrm dx\right)^{1/p}
 \le C_dr
       \left(\frac1{|B|}\int_B|\nabla u|^p\,\mathrm dx\right)^{1/p}.
\end{equation}
在有限 \(p\) 时，这个式子同时比较平均振幅与平均梯度大小。它适用于内部球，不要求全域边界规则，也不包含尚未指定的边界条件。

\ProseParagraphBreak

两个不同半径的均值也可以比较。若 \(B_r=B(x,r)\)、\(B_{2r}=B(x,2r)\subset\Omega\)，则
\[
 \begin{aligned}
 |u_{B_r}-u_{B_{2r}}|
 &\le \frac1{|B_r|}\int_{B_r}|u-u_{B_{2r}}|\\
 &\le C_d r^{1-d/p}\|\nabla u\|_{L^p(B_{2r})}.
 \end{aligned}
\]
第二步先扩大积分区域，再使用 \(L^1\) Poincaré 与 Hölder 不等式。这个半径幂是\secref{b1:sec:2403} Morrey 理论的关键：当 \(p>d\) 时，沿半径逐次减半得到一个可和的几何级数；当 \(p=d\) 时，仅凭这个界不再获得衰减。

\ProseParagraphBreak

另一种局部形式来自零边界估计。对 \(\eta\in C_c^\infty(B_R)\)、\(u\in W^{1,p}(B_R)\)，有限 \(p\) 时 \(\eta u\in W_0^{1,p}(B_R)\)，因此
\[
 \|\eta u\|_{L^p(B_R)}
 \le C_dR\bigl(\|\eta\nabla u\|_{L^p(B_R)}
                   +\|u\nabla\eta\|_{L^p(B_R)}\bigr).
\]
要使右侧的零阶项更小，常把 \(u\) 换成 \(u-u_{B_R}\)，再用均值版本控制它。这种减常数再截断的组织，与上一章 Caccioppoli 不等式中任意常数 \(c\) 的选择一致；局部化引入的截断导数仍必须保留。

% !TeX root = ../../Book1.tex
\section{Sobolev 不等式}
\label{b1:sec:2402}

一阶导数可积，能否使函数本身具有更高的可积指数？Poincaré 型估计控制函数减去常数后的大小，而本节要提高左端的积分指数。证明从紧支集函数开始：函数沿每条坐标线都在两端消失，因而可以用该方向的导数恢复函数值。把各方向的估计合在一起，维数便进入了可积指数。

\ProseParagraphBreak

Nirenberg 的《On elliptic partial differential equations》\cite[Lecture II, pp. 128--129]{Nirenberg1959Elliptic}给出逐坐标积分的证明及幂次代换。本节把其中的乘积积分步骤写成任意维归纳，再用已经建立的稠密性和延拓定理进入 Sobolev 空间。除最后的一维说明外，始终设 \(d\geq2\)；标量域为 \(\mathbb K=\mathbb R\) 或 \(\mathbb C\)，梯度长度采用
\[
 |\nabla u|=\left(\sum_{i=1}^d|\partial_i u|^2\right)^{1/2}.
\]
其 \(L^p\) 范数与一阶导数分量的有限乘积范数等价，所需比较常数只依赖于 \(d,p\)。

% 实读正式来源 Nirenberg1959Elliptic，Numdam 开放扫描：
% https://numdam.org/item/ASNSP_1959_3_13_2_115_0.pdf
% 印刷128--129页（PDF第15--16页），式(2.4)、(2.4)'及三维逐坐标证明。
% 原文一般维数写为同样使用Holder，一般p仅给幂次；本节展开归纳、幂次计算与复值C1细节。
% 一般W1p推广、W0补零和Lipschitz域推广调用本书第22章，不将原文的光滑域讨论扩写成任意域结论。

\subsection{\texorpdfstring{\(1\leq p<d\)}{1≤p<d}}
\label{b1:subsec:240201}

先处理 \(p=1\)。对 \(x=(x_1,\ldots,x_d)\)，以 \(\widehat{x}_i\) 表示删去第 \(i\) 个坐标所得的向量。沿第 \(i\) 条坐标线积分导数后，得到的函数不再依赖 \(x_i\)。下面的不等式正是为这种结构准备的。

\begin{lemma}\label{b1:lem:coordinate-product-integral}
设 \(f_i\in L^1(\mathbb R^{d-1})\) 非负，\(1\leq i\leq d\)。则
\begin{equation}\label{b1:eq:coordinate-product-integral}
 \int_{\mathbb R^d}\prod_{i=1}^d
          f_i(\widehat{x}_i)^{1/(d-1)}\,\mathrm dx
 \leq\prod_{i=1}^d
          \left(\int_{\mathbb R^{d-1}}f_i\right)^{1/(d-1)}.
\end{equation}
\end{lemma}
\begin{proof}
对维数归纳。当 \(d=2\) 时，左端为
\(\int_{\mathbb R^2}f_1(x_2)\cdot f_2(x_1)\,\mathrm dx\)，由 Tonelli 定理恰好等于两个积分的乘积。

设 \(d\geq3\)，且结论在 \(d-1\) 维成立。写 \(x'=(x_1,\ldots,x_{d-1})\)，对 \(i<d\) 置
\[
 F_i(\widehat{x'}_i)=\int_{\mathbb R}f_i(\widehat{x}_i)\,\mathrm dx_d.
\]
它是 \(d-2\) 个变量的非负可积函数，而且
\(\int F_i=\int f_i\)。对固定 \(x'\)，在 \(x_d\) 中使用 \(d-1\) 个指数都为 \(d-1\) 的 Hölder 不等式，得
\[
 \int_{\mathbb R}\prod_{i=1}^{d-1}
       f_i(\widehat{x}_i)^{1/(d-1)}\,\mathrm dx_d
 \leq\prod_{i=1}^{d-1}
       F_i(\widehat{x'}_i)^{1/(d-1)}.
\]
这里的多因子 Hölder 可由两因子形式归纳得到：分出第一个因子，用指数 \(N,N/(N-1)\)，再对后 \(N-1\) 个因子的 \(N/(N-1)\) 次幂应用归纳假设，即得 \(N\) 个指数都为 \(N\) 的形式。

记\eqref{b1:eq:coordinate-product-integral}左端为 \(I\)。由于 \(f_d\) 与 \(x_d\) 无关，上述估计给出
\[
 \begin{aligned}
 I
 &\leq\int_{\mathbb R^{d-1}}f_d(x')^{1/(d-1)}
             \prod_{i=1}^{d-1}F_i(\widehat{x'}_i)^{1/(d-1)}\,\mathrm dx'\\
 &\leq\left(\int_{\mathbb R^{d-1}}f_d\right)^{1/(d-1)}
       \left(\int_{\mathbb R^{d-1}}
             \prod_{i=1}^{d-1}F_i(\widehat{x'}_i)^{1/(d-2)}
             \,\mathrm dx'\right)^{(d-2)/(d-1)}.
 \end{aligned}
\]
第二步的两个 Hölder 指数为 \(d-1\) 和 \((d-1)/(d-2)\)。对最后的积分应用 \(d-1\) 维归纳假设，再取外面的幂，得到
\[
 I\leq\left(\int f_d\right)^{1/(d-1)}
             \prod_{i=1}^{d-1}\left(\int F_i\right)^{1/(d-1)}.
\]
这就是所需估计。全部被积函数非负，交换积分由 Tonelli 定理保证，无须先假定乘积可积。
\end{proof}

把它用于各坐标方向的导数，就得到\term[gagliardo-nirenberg-sobolev-budengshi]{Gagliardo–Nirenberg–Sobolev 不等式}[Gagliardo--Nirenberg--Sobolev inequality]的端点版本。

\begin{theorem}[Gagliardo–Nirenberg–Sobolev，\(p=1\)]\label{b1:thm:gns-endpoint}
若 \(u\in C_c^1(\mathbb R^d;\mathbb K)\)，则
\begin{equation}\label{b1:eq:gns-endpoint}
 \|u\|_{d/(d-1)}
 \leq\frac12\prod_{i=1}^d\|\partial_i u\|_1^{1/d}
 \leq\frac12\|\nabla u\|_1.
\end{equation}
\end{theorem}
\begin{proof}
固定其余坐标。由于 \(u\) 在该坐标线两端为零，微积分基本定理分别给出
\[
 u(x)=\int_{-\infty}^{x_i}\partial_i u(x_1,\ldots,t,\ldots,x_d)\,\mathrm dt
     =-\int_{x_i}^{\infty}\partial_i u(x_1,\ldots,t,\ldots,x_d)\,\mathrm dt.
\]
这对复值函数也成立，只须按实部和虚部积分。对两式分别取绝对值并相加，得
\[
 2|u(x)|\leq
 A_i(\widehat{x}_i),\quad
 A_i(\widehat{x}_i)
 =\int_{\mathbb R}|\partial_i u(x_1,\ldots,t,\ldots,x_d)|\,\mathrm dt.
\]
对 \(i=1,\ldots,d\) 相乘，再取 \(1/(d-1)\) 次幂并积分。由前一引理，
\[
 \int |u|^{d/(d-1)}
 \leq2^{-d/(d-1)}
       \prod_{i=1}^d\left(\int A_i\right)^{1/(d-1)}
 =2^{-d/(d-1)}
       \prod_{i=1}^d\|\partial_i u\|_1^{1/(d-1)}.
\]
取 \((d-1)/d\) 次幂即得第一个不等式。又因每个
\(\|\partial_i u\|_1\leq\|\nabla u\|_1\)，得到第二个不等式。
\end{proof}

这里并未寻求最佳常数。重要的是右端只含一阶导数，而左端指数已经大于一；各坐标方向都参与估计，不能只在一个方向积分后便把其余变量略去。

\subsection{临界指数}
\label{b1:subsec:240202}

\begin{definition}\label{b1:def:sobolev-critical-exponent}
设 \(1\leq p<d\)。定义相应的\term[sobolev-linjie-zhishu]{Sobolev 临界指数}[Sobolev critical exponent]为
\begin{equation}\label{b1:eq:sobolev-critical-exponent}
 p^*\coloneqq\frac{dp}{d-p},\quad
 \frac1{p^*}=\frac1p-\frac1d.
\end{equation}
\end{definition}
\symidx{\(p^*\)：Sobolev 临界指数}

这里的“临界”指给定 \(p<d\) 后，函数值可以达到的边界指数 \(p^*\)，不是把导数指数设成 \(p=d\)。后者将在\secref{b1:sec:2405}另行讨论。比如 \(d=3,p=2\) 时，\(p^*=6\)：平方可积的一阶导数将给出函数的六次可积性。

\ProseParagraphBreak

端点估计应用于函数的一定幂，就能提高导数端的指数。设 \(1<p<d\)，希望把\eqref{b1:eq:gns-endpoint}用于 \(v=|u|^\gamma\)。左端积分中的幂为 \(\gamma d/(d-1)\)，而右端使用 \(p,p'=p/(p-1)\) 的 Hölder 不等式后，会出现 \((\gamma-1)p'\)。要让这两者相同，须取
\begin{equation}\label{b1:eq:sobolev-power-choice}
 \gamma=\frac{p(d-1)}{d-p}>1,\quad
 \frac{\gamma d}{d-1}=(\gamma-1)p'=p^*.
\end{equation}
因此这个幂不是独立的技巧性参数，而是由两端积分必须相容决定的。

\begin{lemma}\label{b1:lem:gns-power-estimate}
设 \(1<p<d\)，\(u\in C_c^1(\mathbb R^d;\mathbb K)\)，并令 \(\gamma\) 如\eqref{b1:eq:sobolev-power-choice}。则
\[
 \|u\|_{p^*}\leq\frac{\gamma}{2}\|\nabla u\|_p.
\]
\end{lemma}
\begin{proof}
先说明幂次代换的正则性。把 \(\mathbb C\) 视为 \(\mathbb R^2\)，映射 \(F(z)=|z|^\gamma\) 在 \(z\ne0\) 处的导数长度为 \(\gamma|z|^{\gamma-1}\)。在原点，
\[
 \frac{|F(z)-F(0)|}{|z|}=|z|^{\gamma-1}\longrightarrow0,
\]
所以导数为零；上述导数长度也趋于零，故 \(F\) 在原点为 \(C^1\)。实值情形是同一论证的一维版本。因此 \(v=|u|^\gamma\) 属于 \(C_c^1\)，不必把它误当成无限次光滑函数。

经典链式法则及欧氏范数不等式给出
\[
 |\nabla v|\leq\gamma|u|^{\gamma-1}\cdot|\nabla u|.
\]
复值情况下，每个偏导数为
\(\partial_i v=\gamma|u|^{\gamma-2}\operatorname{Re}
(\overline u\,\partial_i u)\)（在 \(u\ne0\) 处），在 \(u=0\) 处为零；用
\(\bigl|\operatorname{Re}(\overline u\,\partial_i u)\bigr|
\leq|u|\cdot|\partial_i u|\) 后求平方和，即得同一梯度界。

将 \(v\) 代入端点估计，再用 Hölder 不等式和\eqref{b1:eq:sobolev-power-choice}，得到
\[
 \begin{aligned}
 \|u\|_{p^*}^{\gamma}
 &=\|v\|_{d/(d-1)}
 \leq\frac{\gamma}{2}
          \int |u|^{\gamma-1}\cdot|\nabla u|\\
 &\leq\frac{\gamma}{2}
       \left(\int|u|^{(\gamma-1)p'}\right)^{1/p'}
       \|\nabla u\|_p
 =\frac{\gamma}{2}\|u\|_{p^*}^{\gamma-1}
                  \cdot\|\nabla u\|_p.
 \end{aligned}
\]
紧支集和连续性保证出现的范数有限。若 \(u\ne0\)，约去
\(\|u\|_{p^*}^{\gamma-1}\)；零函数则直接满足结论。
\end{proof}

\subsection{Gagliardo–Nirenberg–Sobolev 不等式}
\label{b1:subsec:240203}

接下来把光滑函数上的估计传给弱导数。需要辨明左端极限：逼近原先只在 \(W^{1,p}\) 中给出，而提高指数后的极限由刚才的估计产生，不能预先假定二者相同。

\begin{theorem}[Gagliardo–Nirenberg–Sobolev]\label{b1:thm:sobolev-subcritical-embedding}
设 \(d\geq2\)、\(1\leq p<d\)。则每个 \(u\in W^{1,p}(\mathbb R^d;\mathbb K)\) 都属于 \(L^{p^*}(\mathbb R^d)\)，并有
\begin{equation}\label{b1:eq:sobolev-global-embedding}
 \|u\|_{p^*}\leq C_{d,p}\|\nabla u\|_p.
\end{equation}
可以取 \(C_{d,1}=1/2\)，当 \(p>1\) 时取
\(C_{d,p}=p(d-1)/[2(d-p)]\)。本节不主张这些常数最佳。
\end{theorem}
\begin{proof}
由\cref{b1:thm:whole-space-sobolev-density}取
\(u_j\in C_c^\infty(\mathbb R^d;\mathbb K)\)，使 \(u_j\to u\) 于 \(W^{1,p}\)。将已经证明的估计用于 \(u_j-u_\ell\)，得
\[
 \|u_j-u_\ell\|_{p^*}
 \leq C_{d,p}\|\nabla u_j-\nabla u_\ell\|_p\longrightarrow0.
\]
由 \(L^{p^*}\) 完备性，存在 \(v\in L^{p^*}\)，使 \(u_j\to v\) 于该空间。取一条子列，使它同时几乎处处收敛到原来的 \(L^p\) 极限 \(u\) 和新的 \(L^{p^*}\) 极限 \(v\)：例如令两种范数误差的相应次幂之和沿该子列可和，再用 Tonelli 定理。因此 \(u=v\) 几乎处处。最后在
\(\|u_j\|_{p^*}\leq C_{d,p}\|\nabla u_j\|_p\) 中取极限，得到结论。这里 \(p^*<\infty\)，所以两种有限指数的逼近都适用。
\end{proof}

\begin{corollary}\label{b1:cor:domain-sobolev-embedding}
设 \(1\leq p<d\)，则有以下两种域上的结论。
\begin{enumerate}
\item\label{b1:item:domain-sobolev-zero-boundary}
对任意非空开集 \(\Omega\subseteq\mathbb R^d\) 及
\(u\in W_0^{1,p}(\Omega)\)，有
\[
 \|u\|_{L^{p^*}(\Omega)}
 \leq C_{d,p}\|\nabla u\|_{L^p(\Omega)}.
\]
常数与 \(\Omega\) 无关。
\item\label{b1:item:domain-sobolev-extension}
若 \(\Omega\) 为有界 Lipschitz 区域，则每个
\(u\in W^{1,p}(\Omega)\) 满足
\[
 \|u\|_{L^{p^*}(\Omega)}
 \leq C_{\Omega,p}\|u\|_{W^{1,p}(\Omega)}.
\]
\end{enumerate}
若相应的 \(\Omega\) 具有有限测度，这两条结论还分别给出所有
\(1\leq q\leq p^*\) 的 \(L^q\) 估计，只须在左端变换指数时加入因子
\(m_d(\Omega)^{1/q-1/p^*}\)。
\end{corollary}
\begin{proof}
第一条使用\cref{b1:prop:w0-zero-extension}：补零函数
\(E_0u\in W^{1,p}(\mathbb R^d)\)，其梯度也由域内梯度补零得到。对它应用\eqref{b1:eq:sobolev-global-embedding}，两端范数分别等于域内范数。

第二条使用\cref{b1:thm:lipschitz-domain-extension}的有界延拓 \(E\)，依次得到
\[
 \|u\|_{L^{p^*}(\Omega)}
 \leq\|Eu\|_{L^{p^*}(\mathbb R^d)}
 \leq C_{d,p}\|\nabla(Eu)\|_{L^p(\mathbb R^d)}
 \leq C_{\Omega,p}\|u\|_{W^{1,p}(\Omega)}.
\]
最后，有限测度下直接对 \(|u|^q\cdot1\) 使用 Hölder 不等式，有
\(\|u\|_q\leq m_d(\Omega)^{1/q-1/p^*}\|u\|_{p^*}\)，即得较低指数的估计。
\end{proof}

一般开集的第一条结论依赖零边界闭包，而非边界光滑性。第二条没有零边界条件，右端因此保留零阶项：非零常数函数在有界区域内属于 \(W^{1,p}\)，梯度却为零，已经排除了对所有函数只用梯度控制的可能。

\begin{proposition}\label{b1:prop:sobolev-subcritical-interpolation}
设 \(u\in W^{1,p}(\mathbb R^d)\)，\(p\leq q\leq p^*\)，并置
\(\theta=d(1/p-1/q)\in[0,1]\)。则
\begin{equation}\label{b1:eq:sobolev-subcritical-interpolation}
 \|u\|_q
 \leq \|u\|_p^{1-\theta}\cdot\|u\|_{p^*}^{\theta}
 \leq C_{d,p}^{\theta}
       \|u\|_p^{1-\theta}\cdot\|\nabla u\|_p^{\theta}.
\end{equation}
端点 \(\theta=0,1\) 按相应单个范数理解。
\end{proposition}
\begin{proof}
只需处理 \(0<\theta<1\)。关系
\[
 \frac1q=\frac{1-\theta}{p}+\frac{\theta}{p^*}
\]
使 \(p/[(1-\theta)q]\) 与 \(p^*/(\theta q)\) 成为共轭指数。将
\(|u|^q=|u|^{(1-\theta)q}\cdot|u|^{\theta q}\) 代入 Hölder 不等式，并取 \(q\) 次根，得到第一个估计；再代入全空间的 Sobolev 不等式，得到第二个。两个端点分别是恒等式和\eqref{b1:eq:sobolev-global-embedding}。
\end{proof}

全空间的测度无限，所以不能像有界域那样，仅凭 \(L^{p^*}\) 就向所有较低指数下降。这里 \(L^p\) 项与梯度项共同控制中间指数；\(q<p\) 不在这个插值结论的范围内。

\subsection{缩放}
\label{b1:subsec:240204}

临界指数还可以从尺度直接读出。固定非零
\(u\in C_c^\infty(\mathbb R^d)\)，令 \(u_\lambda(x)=u(\lambda x)\)，其中 \(\lambda>0\)。换元给出
\[
 \|u_\lambda\|_q=\lambda^{-d/q}\|u\|_q,\quad
 \|\nabla u_\lambda\|_p=\lambda^{1-d/p}\|\nabla u\|_p,
\]
\(q=\infty\) 时按 \(1/q=0\) 理解。如果存在不随函数改变的常数，使所有紧支集光滑函数都满足
\(\|u\|_q\leq C\|\nabla u\|_p\)，那么必有
\[
 \lambda^{d/p-1-d/q}\|u\|_q\leq C\|\nabla u\|_p
 \quad(\lambda>0).
\]
指数为正时令 \(\lambda\to\infty\)，为负时令 \(\lambda\to0\)，都会矛盾。因此必须
\(1/q=1/p-1/d\)，即 \(q=p^*\)。尺度只给出必要条件；前面的坐标积分证明才保证这个指数确实能够达到。

\ProseParagraphBreak

即使定义域固定为一个小球，非齐次估计也不能越过 \(p^*\)。取
\(x_0\in\Omega\)、非零 \(\varphi\in C_c^\infty\bigl(B(0,1)\bigr)\)，并在
\(\varepsilon<\operatorname{dist}(x_0,\Omega^c)\) 时置
\[
 v_\varepsilon(x)
 =\varepsilon^{1-d/p}
       \varphi\left(\frac{x-x_0}{\varepsilon}\right).
\]
这些函数属于 \(C_c^\infty(\Omega)\)，而
\[
 \|v_\varepsilon\|_p=\varepsilon\|\varphi\|_p,\quad
 \|\nabla v_\varepsilon\|_p=\|\nabla\varphi\|_p,\quad
 \|v_\varepsilon\|_q
 =\varepsilon^{1-d/p+d/q}\|\varphi\|_q.
\]
当 \(q>p^*\) 时，最后的指数为负，左端趋于无穷，前两个量却有界。这排除了任何统一的
\(W^{1,p}(\Omega)\rightarrow L^q(\Omega)\) 范数估计；它只检验可积性估计的边界，不需要引入后文的紧性理论。

\ProseParagraphBreak

一维情形须单独读。条件 \(1\leq p<d\) 在 \(d=1\) 时为空，但沿整条实轴的基本定理仍给出另一个端点：
\[
 \|u\|_\infty\leq\frac12\|u'\|_1
 \quad\bigl(u\in W^{1,1}(\mathbb R)\bigr).
\]
对 \(C_c^1\) 函数，它就是从左右两端积分后相加的估计。对一般
\(W^{1,1}\) 函数，取紧支集光滑逼近；将估计用于两项之差，得到一致 Cauchy 列，其连续极限与 \(L^1\) 极限几乎处处相同，故同一界传到极限。这里的结论是本性上确界控制；\secref{b1:sec:2403}讨论 \(p>d\) 时更精确的连续性估计。

% !TeX root = ../../Book1.tex
\section{Morrey 理论}
\label{b1:sec:2403}

当梯度的可积指数超过空间维数时，局部平均振荡不仅趋零，而且在逐次缩小半径后可以求和。这使 Sobolev 函数获得连续代表，并给出两点之间的定量控制。以下先讨论 \(d\ge1\)、\(d<p<\infty\)，令
\[
 \alpha=1-\frac dp\in(0,1).
\]

为把本节放回整章的位置，\cref{b1:tab:first-order-sobolev-regimes}列出 \(d\ge2\) 时一阶理论的三个指数区间。它只给出典型目标；区域、边界条件与端点的精确假设仍以各正式定理为准。
\begin{table}[htbp]
  \centering
  \caption[一阶 Sobolev 的三个区间]{一阶 Sobolev 的三个区间。梯度尺度从积分增益，经临界端点，过渡到 Hölder 连续性。}
  \label{b1:tab:first-order-sobolev-regimes}
  \begin{tabularx}{\textwidth}{@{}p{2.7cm}p{3.6cm}X@{}}
    \toprule
    指数范围 & 临界尺度 & 典型结论 \\
    \midrule
    \(1\le p<d\) & \(p^*=dp/(d-p)\) & 嵌入到 \(L^{p^*}\)，紧性只在严格低于临界指数时出现 \\
    \(p=d\) & \(p^*=\infty\) 的形式端点 & 一般无 \(L^\infty\) 控制；有界域上有全部有限 \(L^q\) 控制，零边界时还有指数可积性 \\
    \(p>d\) & \(\alpha=1-d/p>0\) & 存在 \(C^{0,\alpha}\) 代表，并有 Morrey 型双点估计 \\
    \bottomrule
  \end{tabularx}
\end{table}

以下以实值函数为主，复值或有限维向量值函数可按分量处理。记 \(\|\nabla u\|_{L^p}\) 为弱梯度的欧氏长度的 \(L^p\) 范数；它与本书按偏导分量计算的梯度范数只相差依赖维数及 \(p\) 的等价常数。

Hajłasz 的《Sobolev spaces on metric-measure spaces》\cite[Section 9]{Hajlasz2003SobolevMetric}用球均值的望远镜求和把平均振荡变成双点估计。下面保留这一组织，但利用欧氏体积与 \(p>d\)，直接让球均值在每个点收敛，不先把讨论限制在 Lebesgue 点上。
% 实读 Hajlasz2003SobolevMetric：作者正式公开PDF，§9 Theorem 9.4及完整证明，
% 印刷204--205页；另读Theorem 9.7(3)的Hölder结论陈述，印刷206页。
% 这里改编逐半径球均值比较；一般度量空间的先修与9.7完整证明不作为本节依赖。
% 欧氏球Poincare由本章24.1提供；每一点的极限、任意半径一致性及代表辨析另行展开。

\subsection{\texorpdfstring{\(p>d\)}{p>d}}
\label{b1:subsec:240301}

对球 \(B=B(x,r)\)，记 \(u_B=m_d(B)^{-1}\int_Bu\)。球上的 Poincaré 不等式与 Hölder 不等式给出
\begin{equation}\label{b1:eq:morrey-average-oscillation}
 \begin{split}
 \frac1{m_d(B)}\int_B|u-u_B|
 &\le m_d(B)^{-1/p}\|u-u_B\|_{L^p(B)}\\
 &\le C_{d,p}r^{1-d/p}\|\nabla u\|_{L^p(B)}.
 \end{split}
\end{equation}
这里使用的是\cref{b1:thm:ball-poincare-lp}，无需先知道 \(u\) 连续。幂次 \(1-d/p\) 是导数贡献的一个长度因子，与平均化引入的 \(r^{-d/p}\) 相乘所得。当它为正时，各个尺度的误差形成收敛的几何级数。

\begin{lemma}\label{b1:lem:morrey-ball-average-limit}
若 \(u\in W^{1,p}_{\mathrm{loc}}(\Omega)\)，其中 \(\Omega\subset\mathbb R^d\) 为开集、\(d<p<\infty\)，则对每个 \(x\in\Omega\)，有限极限
\[
 u^*(x)=\lim_{r\downarrow0}u_{B(x,r)}
\]
都存在，且 \(u^*=u\) 几乎处处。对每个满足 \(\overline{B(x,R)}\Subset\Omega\) 的球，还有
\begin{equation}\label{b1:eq:morrey-average-error}
 |u^*(x)-u_{B(x,R)}|
 \le C_{d,p}R^\alpha\|\nabla u\|_{L^p(B(x,R))}.
\end{equation}
\end{lemma}
\begin{proof}
固定 \(x,R\)，置 \(r_j=2^{-j}R\) 和 \(B_j=B(x,r_j)\)。因为两相邻球的体积比为 \(2^d\)，
\[
 \begin{split}
 |u_{B_{j+1}}-u_{B_j}|
 &\le\frac1{m_d(B_{j+1})}\int_{B_j}|u-u_{B_j}|\\
 &\le C_{d,p}r_j^\alpha\|\nabla u\|_{L^p(B_j)}
 \le C_{d,p}2^{-j\alpha}R^\alpha
                         \|\nabla u\|_{L^p(B_0)}.
 \end{split}
\]
右端可求和，所以均值序列为 Cauchy 列，存在有限极限；从 \(j=0\) 起求和就得到\eqref{b1:eq:morrey-average-error}。

还要排除极限对选定半径列的依赖。若 \(r_{j+1}<r\le r_j\)，则 \(B(x,r)\subset B_j\)，且体积比仍不超过 \(2^d\)。同一估计给出
\[
 |u_{B(x,r)}-u_{B_j}|
 \le C_{d,p}r_j^\alpha\|\nabla u\|_{L^p(B_0)}
 \longrightarrow0.
\]
因此所有趋零半径的均值都趋于同一极限，它与最初的 \(R\) 无关。以上论证对每个固定点成立，没有删除例外点。最后由\cref{b1:thm:lebesgue-point}，在几乎每个点，该极限等于原函数值，故确实得到 \(u\) 的一个代表。
\end{proof}

当 \(p=d\) 时，上面用整个固定球上的梯度范数估计的几何级数失去衰减。这个观察只说明当前证明不能给出正的 Hölder 指数。\secref{b1:sec:2405}将讨论临界情形中的其他控制。

\subsection{全空间 Morrey 不等式}
\label{b1:subsec:240302}

\begin{definition}\label{b1:def:holder-space}
给定非空集合 \(A\subset\mathbb R^d\) 及 \(0<\gamma\le1\)，若函数 \(v:A\rightarrow\mathbb R\) 满足
\[
 [v]_{0,\gamma;A}
 \coloneqq\sup_{\substack{x,y\in A\\x\ne y}}
       \frac{|v(x)-v(y)|}{|x-y|^\gamma}<\infty,
\]
则称它在 \(A\) 上为 \(\gamma\) 阶\term[holder-lianxu]{Hölder 连续}[Hölder continuous]。若还满足 \(\sup_A|v|<\infty\)，记 \(v\in C^{0,\gamma}(A)\)，并规定
\[
 \|v\|_{C^{0,\gamma}(A)}
 \coloneqq\sup_A|v|+[v]_{0,\gamma;A}.
\]
集合只有一个点时，半范数中的上确界约定为零。
\end{definition}
\symidx{\(C^{0,\gamma}(A)\)：有界Hölder连续函数空间}
\symidx{\([v]_{0,\gamma;A}\)：Hölder半范数}

指数为 \(1\) 时就是 Lipschitz 条件。对紧集，有限半范数和任意一个有限点值已保证有界；在全空间上则不能省去零阶条件，例如 \(v(x)=x_1\) 虽然 Lipschitz，却不属于这里的 \(C^{0,1}(\mathbb R^d)\)。

\begin{theorem}[Morrey]\label{b1:thm:morrey-whole-space}
设 \(d<p<\infty\) 且 \(u\in W^{1,p}(\mathbb R^d)\)。球平均极限 \(u^*\) 是唯一的连续代表，属于 \(C^{0,\alpha}(\mathbb R^d)\)，其中 \(\alpha=1-d/p\)，并满足
\begin{equation}\label{b1:eq:morrey-whole-space}
 [u^*]_{0,\alpha;\mathbb R^d}
 \le C_{d,p}\|\nabla u\|_{L^p(\mathbb R^d)},
 \quad
 \|u^*\|_{C^{0,\alpha}(\mathbb R^d)}
 \le C_{d,p}\|u\|_{W^{1,p}(\mathbb R^d)}.
\end{equation}
\end{theorem}
\begin{proof}
给定 \(x\ne y\)，令 \(r=|x-y|\)，并取
\[
 B_x=B(x,r),\quad B_y=B(y,r),\quad B_*=B(x,3r).
\]
两个小球都包含于 \(B_*\)，体积比只依赖于 \(d\)。由\eqref{b1:eq:morrey-average-oscillation}，
\[
 |u_{B_x}-u_{B_*}|+|u_{B_y}-u_{B_*}|
 \le C_{d,p}r^\alpha\|\nabla u\|_{L^p(B_*)}.
\]
再在 \(x,y\) 两点使用\eqref{b1:eq:morrey-average-error}，得到
\begin{equation}\label{b1:eq:morrey-local-pair}
 |u^*(x)-u^*(y)|
 \le C_{d,p}|x-y|^\alpha\cdot
                 \|\nabla u\|_{L^p(B(x,3|x-y|))}.
\end{equation}
放大积分域到整个空间，即得 Hölder 半范数界，也证明 \(u^*\) 处处连续。

半范数不能单独控制常数，因此还须估计一个平均值。对任意 \(x\)，用半径为 \(1\) 的球，
\[
 \begin{split}
 |u^*(x)|
 &\le|u^*(x)-u_{B(x,1)}|+|u_{B(x,1)}|\\
 &\le C_{d,p}\|\nabla u\|_{L^p(\mathbb R^d)}
       +\omega_d^{-1/p}\|u\|_{L^p(\mathbb R^d)}.
 \end{split}
\]
于是得到全范数估计。若另一个连续函数与 \(u\) 几乎处处相等，它与 \(u^*\) 的差连续且几乎处处为零；一旦在某点非零，就会在一个正测度小球内非零，矛盾。因此连续代表唯一。
\end{proof}

上述\term[morrey-budengshi]{Morrey 不等式}[Morrey inequality]把弱梯度的积分控制转成点值的变化控制。证明中的每一个平均值只依赖等价类，连续代表不是通过任意指定某个零测集上的值获得的。

\subsection{局部与区域版本}
\label{b1:subsec:240303}

全空间估计也给出内部估计。这里必须留意参与均值比较的球是否仍在原开集内，而不能直接对靠近边界的点使用越出开集的球。

\begin{corollary}\label{b1:cor:morrey-domain}
设 \(d<p<\infty\)、\(\alpha=1-d/p\)。
若 \(u\in W^{1,p}_{\mathrm{loc}}(\Omega)\)，则其球平均代表局部 Hölder 连续。具体地，若 \(\overline{B(a,8R)}\Subset\Omega\)，则
\begin{equation}\label{b1:eq:morrey-interior-ball}
 \begin{split}
 [u^*]_{0,\alpha;\overline{B(a,R)}}
 &\le C_{d,p}\|\nabla u\|_{L^p(B(a,8R))},\\
 \sup_{\overline{B(a,R)}}|u^*|
 &\le C_{d,p}\left(
        R^{-d/p}\|u\|_{L^p(B(a,8R))}
        +R^\alpha\|\nabla u\|_{L^p(B(a,8R))}\right).
 \end{split}
\end{equation}
若 \(\overline U\Subset V\subset\Omega\)，其中 \(U,V\) 为开集，且 \(u\in W^{1,p}(V)\)，则
\[
 \|u^*\|_{C^{0,\alpha}(\overline U)}
 \le C_{U,V,d,p}\|u\|_{W^{1,p}(V)}.
\]
若 \(\Omega\) 为有界 Lipschitz 区域，则 \(u\in W^{1,p}(\Omega)\) 有唯一的连续延伸到 \(\overline\Omega\) 的代表，且
\begin{equation}\label{b1:eq:morrey-lipschitz-domain}
 \|u^*\|_{C^{0,\alpha}(\overline\Omega)}
 \le C_{\Omega,p}\|u\|_{W^{1,p}(\Omega)}.
\end{equation}
\end{corollary}
\begin{proof}
对 \(x,y\in\overline{B(a,R)}\)，有 \(|x-y|\le2R\)，故
\(B(x,3|x-y|)\subset B(a,8R)\)。主定理的双点比较只在这些内部球中进行，所以\eqref{b1:eq:morrey-local-pair}仍成立，给出第一个估计。对 \(x\) 使用 \(B(x,R)\) 的平均值，并以 Hölder 不等式控制该平均值，就得到第二个估计。这同时证明代表的局部 Hölder 连续性。

给定 \(\overline U\Subset V\)，取固定 \(\chi\in C_c^\infty(V)\)，使它在 \(\overline U\) 的一个邻域内等于 \(1\)。将 \(\chi u\) 补零到整个空间。由弱乘积公式及\cref{b1:lem:compact-sobolev-zero-extension}，所得 \(w\) 满足
\[
 \|w\|_{W^{1,p}(\mathbb R^d)}
 \le C_{\chi,p}\|u\|_{W^{1,p}(V)}.
\]
在 \(\overline U\) 附近的充分小球上，\(u\) 与 \(w\) 的平均值相同，故两者的精确代表也相同。应用全空间估计即得第二项结论。

最后，若 \(\Omega\) 有界且为 Lipschitz 区域，用\cref{b1:thm:lipschitz-domain-extension}取得全空间延拓 \(Eu\)，再把其 Hölder 连续代表限制到 \(\overline\Omega\)。在内部，它与原来的球平均代表一致；边界值由内部连续趋近唯一确定。因此结果不依赖所选延拓算子，且延拓范数界给出\eqref{b1:eq:morrey-lipschitz-domain}。
\end{proof}

区域不必连通，但区域版本使用的是完整 \(W^{1,p}\) 范数，不能一律只保留 \(\|\nabla u\|_p\)。例如在两个分离的球上分别取常数 \(0\) 与 \(1\)，弱梯度为零，跨两个分量的 Hölder 半范数却非零。全空间延拓同时处理各分量之间的变化，相关常数包含区域的几何信息。

当 \(p=\infty\) 时，结论中的指数改为 \(1\)，但证明不通过把 \(p\) 代入有限指数公式来完成。\cref{b1:thm:w1infty-lipschitz}已经给出局部 Lipschitz 代表；在全空间中，把连接两点的线段分为有限个落在局部球内的小段并求和，就得到由全域弱梯度本性界控制的 Lipschitz 常数。其上确界由 \(\|u\|_\infty\) 控制。再使用截断或有界 Lipschitz 区域的延拓算子，分别得到局部和闭区域上的 \(C^{0,1}\) 估计。

指数 \(\alpha=1-d/p\) 在统一范数估计中不能提高。固定非恒定函数 \(\phi\in C_c^\infty(B(0,1))\)，并令
\[
 u_\varepsilon(x)=\varepsilon^\alpha\phi(x/\varepsilon),
 \quad 0<\varepsilon<1.
\]
换元直接给出
\[
 \|u_\varepsilon\|_p=\varepsilon\|\phi\|_p,\quad
 \|\nabla u_\varepsilon\|_p=\|\nabla\phi\|_p.
\]
这些函数在 \(W^{1,p}\) 中一致有界。然而，选取 \(\phi(a)\ne\phi(b)\)，对每个 \(\alpha<\beta\le1\)，
\[
 [u_\varepsilon]_{0,\beta;\mathbb R^d}
 \ge\varepsilon^{\alpha-\beta}
       \frac{|\phi(a)-\phi(b)|}{|a-b|^\beta}
 \longrightarrow\infty.
\]
所以不存在把所有 \(W^{1,p}\) 函数的 \(C^{0,\beta}\) 范数控制住的统一常数；同一例子也适用于包含这些小支集的固定球。每个 \(u_\varepsilon\) 本身仍然光滑，这里排除的是更强的连续嵌入估计，而不是说单个函数不能具有更好的正则性。

\subsection{精确代表}
\label{b1:subsec:240304}

\secref{b1:sec:1404}曾由球平均定义精确代表，并在极限不存在处另作赋值。现在，超临界可积性使这一例外情形在内部完全消失。

\begin{proposition}\label{b1:prop:morrey-precise-representative}
若 \(u\in W^{1,p}_{\mathrm{loc}}(\Omega)\)、\(d<p<\infty\)，则其精确代表 \(u^*\) 是唯一的连续代表，并且对每个 \(x\in\Omega\)，
\[
 \lim_{r\downarrow0}\frac1{m_d(B(x,r))}
       \int_{B(x,r)}|u(y)-u^*(x)|\,\mathrm dy=0.
\]
特别地，\(u^*\) 的每个内部点都是 Lebesgue 点。若 \(u_j\to u\) 于 \(W^{1,p}_{\mathrm{loc}}(\Omega)\)，则对每个相对紧开集 \(U\)，精确代表在 \(C^{0,\alpha}(\overline U)\) 中收敛；这里取 \(\overline U\Subset\Omega\)，且 \(\alpha=1-d/p\)。
\end{proposition}
\begin{proof}
有限极限已由\cref{b1:lem:morrey-ball-average-limit}在每个点构造，局部 Hölder 连续性由前一推论给出。连续代表的唯一性仍由非空开球具有正测度得到。对充分小的 \(r\)，
\[
 \begin{split}
 \frac1{m_d(B(x,r))}\int_{B(x,r)}|u-u^*(x)|
 &\le\frac1{m_d(B(x,r))}\int_{B(x,r)}|u-u_{B(x,r)}|
       +|u_{B(x,r)}-u^*(x)|\\
 &\le C_{d,p}r^\alpha
                 \|\nabla u\|_{L^p(B(x,r))}
 \longrightarrow0.
 \end{split}
\]
将积分中的 \(u\) 换成与它几乎处处相等的 \(u^*\)，即得每点的 Lebesgue 点性质。

最后取 \(\overline U\Subset V\) 且 \(\overline V\Subset\Omega\)。球均值的线性及每点极限说明
\((u_j-u)^*=u_j^*-u^*\)。对差函数使用前一推论，便有
\[
 \|u_j^*-u^*\|_{C^{0,\alpha}(\overline U)}
 \le C_{U,V,d,p}\|u_j-u\|_{W^{1,p}(V)}
 \longrightarrow0.
\]
\end{proof}

连续性属于这个被确定的代表，而不是等价类中的每个可测函数。即使只在一点改变 \(u^*\) 的值，也不改变任何 \(L^p\) 范数或弱导数，却可能破坏该点的连续性。另一方面，在有界 Lipschitz 区域上，\cref{b1:thm:lipschitz-lp-trace}说明这个闭域连续代表的边界限制恰是前章的迹；积分方式定义的边界值与通常点值在这里相遇。

本节得到的是函数值的 Hölder 控制，并不使弱梯度自动连续。把同一方法应用到哪些低阶导数、怎样组织高阶连续代表，是\secref{b1:sec:2404}高阶嵌入要处理的问题。

% !TeX root = ../../Book1.tex
\section{高阶 Sobolev 嵌入}
\label{b1:sec:2404}

一阶嵌入可以作用于函数，也可以作用于它的低阶弱导数。逐次使用这一事实，就能把高阶可积性转化为更高的积分指数，或转化为经典导数的 Hölder 连续性。整个过程有两个需要区分的步骤：先提高每个弱导数的可积性，再证明所得连续代表确实互为经典导数。

\ProseParagraphBreak

本节首先取整数 \(k\ge1\) 和有限指数 \(1\le p<\infty\)，在全空间中讨论。一般开集上的内部结论可通过截断取得；涉及整个边界时，所需延拓必须控制到 \(k\) 阶。最后单独从双积分定义证明 \(0<s<1\) 的分数阶版本。
% 整数阶迭代由本章Gagliardo--Nirenberg--Sobolev与Morrey定理自足推出。
% 分数阶均值比较沿用已实读Hajlasz2003SobolevMetric §9 Theorem 9.4
% 的望远镜组织（印刷204--205页）；原文9.7(3)仅作为已读的范围说明，
% 不把该文未展开于本书的度量空间理论当作高阶或分数阶证明前置。

\subsection{高阶导数}
\label{b1:subsec:240401}

若 \(u\in W^{k,p}(\mathbb R^d)\)，则每个 \(|\beta|\le k-1\) 的弱导数 \(\partial^\beta u\) 都属于 \(W^{1,p}\)。在 \(p<d\) 时，对这些函数分别使用一阶嵌入，便得到
\[
 W^{k,p}(\mathbb R^d)\subset W^{k-1,p_1}(\mathbb R^d),
 \quad \frac1{p_1}=\frac1p-\frac1d.
\]
右端包含所有不超过 \(k-1\) 阶的弱导数，不只是最高阶的某几个分量。因此下一次迭代仍有完整的 Sobolev 空间可用。

\begin{theorem}\label{b1:thm:higher-order-subcritical-embedding}
设 \(1\le p<\infty\)、\(k\ge1\) 且 \(kp<d\)，并令
\[
 q=\frac{dp}{d-kp}.
\]
则有连续嵌入
\[
 W^{k,p}(\mathbb R^d)\subset L^q(\mathbb R^d),
 \quad
 \|u\|_q\le C_{d,k,p}\|u\|_{W^{k,p}(\mathbb R^d)}.
\]
\end{theorem}
\begin{proof}
规定 \(p_0=p\) 及
\[
 \frac1{p_j}=\frac1p-\frac jd,\quad 0\le j\le k.
\]
条件 \(kp<d\) 保证这些指数有限，而且每一步开始时 \(p_j<d\)，其中 \(0\le j<k\)。若已知 \(u\in W^{k-j,p_j}\)，则对每个 \(|\beta|\le k-j-1\)，由\cref{b1:thm:sobolev-subcritical-embedding}，
\[
 \|\partial^\beta u\|_{p_{j+1}}
 \le C_{d,p_j}\sum_{i=1}^d
                  \|\partial^{\beta+e_i}u\|_{p_j}.
\]
这些弱导数仍是原来的分布导数；改变可积指数不改变其定义。对有限多个 \(\beta\) 求和，得到
\[
 \|u\|_{W^{k-j-1,p_{j+1}}}
 \le C_{d,k,p_j}\|u\|_{W^{k-j,p_j}}.
\]
从 \(j=0\) 迭代到 \(j=k-1\)，便到达 \(W^{0,p_k}=L^q\)，并得到范数界。
\end{proof}

每提高一次积分指数，消耗一阶导数，同时使 \(1/p\) 减少 \(1/d\)。例如在五维空间中，\(W^{2,2}\) 先嵌入 \(W^{1,10/3}\)，再嵌入 \(L^{10}\)。若中途指数达到 \(d\)，就不能继续把分母等于零的表达式当作下一步的一阶公式；若指数超过 \(d\)，则应转入 Morrey 理论。

\subsection{Hölder 空间}
\label{b1:subsec:240402}

\begin{definition}\label{b1:def:higher-holder-space}
给定开集 \(U\subset\mathbb R^d\)、整数 \(m\ge0\) 及 \(0<\alpha\le1\)，若 \(v\in C^m(U)\)，所有不超过 \(m\) 阶的经典偏导数均有界，且每个 \(m\) 阶偏导数都为 \(\alpha\) 阶 Hölder 连续，则记 \(v\in C^{m,\alpha}(U)\)，并规定
\[
 \|v\|_{C^{m,\alpha}(U)}
 \coloneqq\sum_{|\beta|\le m}\sup_U|\partial^\beta v|
   +\sum_{|\beta|=m}[\partial^\beta v]_{0,\alpha;U}.
\]
记号 \(C^{m,\alpha}(\overline U)\) 还要求这些导数连续延伸到 \(\overline U\)，上式中的上确界与 Hölder 半范数也在闭包上计算。当 \(m=0\) 时与前节的定义一致。
\end{definition}
\symidx{\(C^{m,\alpha}(U)\)：最高\(m\)阶经典导数为Hölder连续的空间}

\begin{theorem}\label{b1:thm:higher-order-holder-embedding}
设 \(k\ge1\)、\(1\le p<\infty\)，且 \(k-d/p>0\) 不是整数。写成
\[
 k-\frac dp=m+\alpha,\quad m\in\mathbb N_0,\quad 0<\alpha<1.
\]
则每个 \(u\in W^{k,p}(\mathbb R^d)\) 都有唯一的 \(C^m\) 代表 \(v\)，满足
\[
 v\in C^{m,\alpha}(\mathbb R^d),\quad
 \|v\|_{C^{m,\alpha}(\mathbb R^d)}
 \le C_{d,k,p}\|u\|_{W^{k,p}(\mathbb R^d)}.
\]
\end{theorem}
\begin{proof}
令 \(\ell=k-m-1\ge0\)，并规定
\[
 \frac1q=\frac1p-\frac\ell d=\frac{1-\alpha}{d}.
\]
故 \(q>d\)。若 \(\ell=0\)，已知 \(u\in W^{m+1,q}\)。若 \(\ell>0\)，前一证明中的一阶迭代可以进行 \(\ell\) 次：最后一次开始前的指数仍小于 \(d\)，因为对应的 \(d/p_{\ell-1}=2-\alpha>1\)。因此两种情形均得到
\[
 \|u\|_{W^{m+1,q}(\mathbb R^d)}
 \le C_{d,k,p}\|u\|_{W^{k,p}(\mathbb R^d)}.
\]
对每个 \(|\beta|\le m\)，函数 \(\partial^\beta u\) 属于 \(W^{1,q}\)。由\cref{b1:thm:morrey-whole-space}，它有 \(C^{0,\alpha}\) 代表 \(v_\beta\)，且相应范数受上式控制。

现在核实这些代表确实是同一经典函数的导数。取局部矩形及其稍大的内部邻域，令 \(u_\varepsilon=\rho_\varepsilon*u\)。由\cref{b1:prop:sobolev-friedrichs-smoothing}，
\[
 \partial^\beta u_\varepsilon
 =\rho_\varepsilon*\partial^\beta u
 =\rho_\varepsilon*v_\beta,\quad |\beta|\le m.
\]
因为 \(v_\beta\) 连续，右端在每个紧集上一致趋于 \(v_\beta\)。对 \(|\beta|<m\)，在矩形内的坐标线段上有
\[
 \partial^\beta u_\varepsilon(x+te_i)-\partial^\beta u_\varepsilon(x)
 =\int_0^t\partial^{\beta+e_i}u_\varepsilon(x+re_i)\,\mathrm dr.
\]
让 \(\varepsilon\) 趋零，得到相同恒等式，其两侧分别换成 \(v_\beta\) 与 \(v_{\beta+e_i}\)。由于被积函数连续，\(v_\beta\) 的经典偏导数就是 \(v_{\beta+e_i}\)。连续偏导保证经典可微；逐阶使用便得到
\(v_0\in C^m\) 且 \(\partial^\beta v_0=v_\beta\)。当 \(m=0\) 时，只需已经得到的连续代表。

对所有低阶上确界及最高阶 Hölder 半范数求和，得到所需估计。连续代表唯一，因而 \(C^m\) 代表也唯一。
\end{proof}

正整数差不能不加论证地解释为端点 Lipschitz 正则性。下面给出足以覆盖任意较小 Hölder 指数的结论。

\begin{proposition}\label{b1:prop:integer-gap-holder-embedding}
设 \(1<p<\infty\)，且 \(j=k-d/p\) 为正整数。对每个 \(0<\alpha<1\)，均有连续嵌入
\[
 W^{k,p}(\mathbb R^d)\subset C^{j-1,\alpha}(\mathbb R^d),
\]
其中右端取经典代表。
\end{proposition}
\begin{proof}
选取充分接近 \(p\) 的 \(p_0\in(1,p)\)，使
\[
 k-\frac d{p_0}=j-1+\alpha_0,\quad \alpha<\alpha_0<1.
\]
固定 \(\chi\in C_c^\infty\bigl(B(0,2)\bigr)\)，使它在 \(B(0,1)\) 上等于 \(1\)，并令
\(\chi_a(x)=\chi(x-a)\)。由于支集体积固定，Hölder 不等式及高阶弱 Leibniz 公式给出与中心 \(a\) 无关的界
\[
 \|\chi_a u\|_{W^{k,p_0}(\mathbb R^d)}
 \le C_{d,k,p,p_0,\chi}\|u\|_{W^{k,p}(\mathbb R^d)}.
\]
应用前一定理，\(\chi_a u\) 在 \(B(a,1)\) 上的经典代表具有一致的 \(C^{j-1,\alpha_0}\) 界。不同球上的代表在交集几乎处处相等，因连续而处处相等，其经典导数也相同，所以这些局部代表拼成同一个 \(C^{j-1}\) 函数。

所有不超过 \(j-1\) 阶的导数由局部界得到全域上确界。对最高阶导数的两点差，若 \(|x-y|<1\)，取中心 \(a=x\)，局部估计与
\(|x-y|^{\alpha_0}\le|x-y|^\alpha\) 给出所需界；若 \(|x-y|\ge1\)，则用两倍上确界控制差值。两种距离范围合起来得到全域的 \(C^{j-1,\alpha}\) 范数界。
\end{proof}

这个证明中的降指数只在固定大小的支集上进行，没有使用全空间 \(L^p\subset L^{p_0}\) 这一错误包含关系。它也没有证明 \(\alpha=1\) 的结论。若 \(p=1\)，就不能再降到允许范围内的 \(p_0<p\)；该临界端点需要专门论证，本节不从上述方法推断其最强结论。差 \(k-d/p=0\) 同样不在这个正整数差命题中。无穷指数若另作讨论，应逐阶使用 Lipschitz 代表的判据，而不是代入有限指数迭代。

\ProseParagraphBreak

一般开集上的内部版本不需要边界正则性：若
\(\overline U\Subset V\)、\(u\in W^{k,p}(V)\)，选取在 \(\overline U\) 附近等于 \(1\) 的 \(\chi\in C_c^\infty(V)\)，以高阶弱乘积公式估计 \(\chi\cdot u\)，再用\cref{b1:lem:compact-sobolev-zero-extension}补零即可应用上述全空间结论。由此取得内部的 \(L^q\) 或 \(C^{m,\alpha}(\overline U)\) 估计。

\ProseParagraphBreak

若要把这些估计延至整个 \(\overline\Omega\)，可以假设存在有界线性延拓
\[
 E_k:W^{k,p}(\Omega)\longrightarrow W^{k,p}(\mathbb R^d),
 \quad (E_ku)|_\Omega=u,
\]
再应用全空间定理。\secref{b1:sec:2203}只为 Lipschitz 区域证明了一阶延拓；在这里讨论 \(k>1\) 时，不能把那个一阶算子当作已经具备高阶有界性。

\subsection{整数阶与非整数阶}
\label{b1:subsec:240403}

整数阶迭代使用了弱导数链。对于 \(0<s<1\) 的空间，正则性由差分双积分度量，必须重新从该定义取得平均振荡估计。两种推导最后会汇合到相同的球均值方法。

\begin{theorem}[分数阶 Morrey]\label{b1:thm:fractional-morrey}
设 \(0<s<1\)、\(1\le p<\infty\) 且 \(sp>d\)，令
\(\alpha=s-d/p\in(0,1)\)。每个 \(u\in W^{s,p}(\mathbb R^d)\) 都有唯一的连续代表 \(u^*\)，其值在每个点由球平均极限给出，且
\[
 [u^*]_{0,\alpha;\mathbb R^d}\le C_{d,s,p}[u]_{s,p},
 \quad
 \|u^*\|_{C^{0,\alpha}(\mathbb R^d)}
 \le C_{d,s,p}\|u\|_{W^{s,p}(\mathbb R^d)}.
\]
\end{theorem}
\begin{proof}
先取 \(B=B(x,r)\)。由平均值的 Jensen 不等式，
\[
 \begin{split}
 \frac1{m_d(B)}\int_B|u-u_B|^p
 &\le\frac1{m_d(B)^2}\int_B\int_B|u(z)-u(y)|^p\,\mathrm dy\,\mathrm dz\\
 &\le\frac{(2r)^{d+sp}}{m_d(B)^2}
       \int_B\int_B
          \frac{|u(z)-u(y)|^p}{|z-y|^{d+sp}}\,\mathrm dy\,\mathrm dz\\
 &\le C_{d,s,p}r^{sp-d}[u]_{s,p;B}^p.
 \end{split}
\]
再用 Hölder 不等式，得到
\[
 \frac1{m_d(B)}\int_B|u-u_B|
 \le C_{d,s,p}r^\alpha[u]_{s,p;B}.
\]
固定任意 \(x,R\)，取 \(B_i=B(x,2^{-i}R)\)。与相邻球的体积比较给出
\[
 |u_{B_{i+1}}-u_{B_i}|
 \le C_{d,s,p}2^{-i\alpha}R^\alpha[u]_{s,p;B(x,R)}.
\]
因为 \(\alpha>0\)，级数收敛。对夹在两个相邻二分半径之间的任意半径，仍可用体积比及刚才的平均振荡界比较均值。因此球平均在每个点存在有限极限，而且
\[
 |u^*(x)-u_{B(x,R)}|
 \le C_{d,s,p}R^\alpha[u]_{s,p;B(x,R)}.
\]
这正是\cref{b1:lem:morrey-ball-average-limit}证明中的求和与任意半径夹逼步骤；这里只将弱梯度范数换成了由双积分直接得到的控制量。Lebesgue 点定理保证该极限几乎处处等于 \(u\)。

给定 \(x\ne y\)，令 \(r=|x-y|\)，比较 \(B(x,r)\)、\(B(y,r)\) 与 \(B(x,3r)\) 的均值。两个小球与大球的体积比有界，所以它们的均值差由大球的平均振荡控制。再加上两个点各自的均值误差，得到
\[
 |u^*(x)-u^*(y)|
 \le C_{d,s,p}|x-y|^\alpha\cdot[u]_{s,p;B(x,3|x-y|)}
 \le C_{d,s,p}|x-y|^\alpha\cdot[u]_{s,p}.
\]
于是 \(u^*\) 连续，并取得所需 Hölder 半范数界。取 \(R=1\)，以 \(\|u\|_p\) 控制单位球均值，便得到全范数界。连续代表的唯一性与前节相同。
\end{proof}

上述论证仍在每个点构造极限，而不只是在几乎处处的点获得一个差分估计。它也说明分数阶嵌入的尺度来自 \(sp-d>0\)：双积分能量控制平均振荡，再由几何级数恢复点值。这里没有通过把整数阶公式中的 \(k\) 直接替换为 \(s\) 来定义分数阶导数。

\ProseParagraphBreak

Hajłasz 的《Sobolev spaces on metric-measure spaces》\cite[Theorem 9.4]{Hajlasz2003SobolevMetric}提供了这种均值比较的系统组织。整数阶和分数阶条件如何产生平均振荡界，各有自己的证明；将这些条件放在同一尺度观点下理解，也有助于辨清\secref{b1:sec:2405}临界情形中哪些步骤不再成立。

高阶嵌入应同时核对指数余量与区域的延拓性质，见\cref{b1:tab:visual-t13}。
% Source fingerprint: b498fd314c29f9c0aadfeeb4aca942a02d5a32b6174e8cf3848275fbd1ed6b88
% T13: raw TeX cells preserved; no numeric highlighting.
\begin{table}[htbp]
\centering\small\renewcommand{\arraystretch}{1.18}\setlength{\tabcolsep}{4pt}
\caption[高阶 Sobolev 嵌入速查]{高阶 Sobolev 嵌入速查。本表补充前面的一阶指数表。函数空间的范数与代表约定均沿用本节，不把内部正则性写成跨边界结论。}
\label{b1:tab:visual-t13}
\begin{tabularx}{\linewidth}{@{}>{\raggedright\arraybackslash}p{3.1cm}>{\raggedright\arraybackslash}p{4.7cm}>{\raggedright\arraybackslash}X@{}}
\toprule
参数区间 & 本节得到的目标 & 适用说明 \\
\midrule
\(kp<d\)，\(1\le p<\infty\) & \(L^q\)，\(q=dp/(d-kp)\) & 逐次一阶嵌入；分母必须严格为正 \\
\(kp=d\) & 不能把上行公式代入为无穷指数 & 需按相应临界理论处理，不由形式代入宣称有界连续 \\
\(k-d/p=m+\alpha\)，\(0<\alpha<1\) & 具有 \(C^{m,\alpha}\) 代表 & 有限 p；余量非整数，经典导数的相容性由正文证明 \\
\(j=k-d/p\in\mathbb N_{>0}\) & \(C^{j-1,\alpha}\)，任意 \(0<\alpha<1\) & 本节该命题要求 \(1<p<\infty\)；不擅取 \(\alpha=1\) \\
一般区域 & 内部结论通过局部化取得；全域结论需相应高阶延拓 & 一阶延拓不能直接替代 \(k\) 阶延拓假设 \\
\bottomrule
\end{tabularx}
\end{table}


% !TeX root = ../../Book1.tex
\OptionalSection{一阶临界嵌入：\texorpdfstring{\(W^{1,d}\)}{W1,d}}{b1:sec:2405}

本节只处理一阶临界情形 \(W^{1,d}\)，不讨论一般 \(W^{k,p}\) 在 \(kp=d\) 时的端点。此时 Sobolev 次临界指数的分母为零，而 Morrey 指数也恰好降为零。这并不意味着所有估计都消失了；它意味着不能继续用前面的幂次公式描述端点。以下设 \(d\ge2\)，先证明零边界函数在每个有限 \(L^q\) 中都有控制，再追踪常数随 \(q\) 的增长，把这些估计合成为指数可积性。

\ProseParagraphBreak

Moser 的《A Sharp Form of an Inequality by N. Trudinger》\cite[1. Statements, p. 1077]{Moser1971Trudinger}说明了临界指数可积性的背景，并进一步研究最佳常数。本节采用较直接的卷积与指数级数路线，只证明足以展示临界现象的非最优常数版本；不把下述常数认作原文的尖锐常数。
% 实读范围：Moser1971Trudinger官方首页PDF，印刷1077页第1节。
% 完整Theorem1跨至1078页，本次未读全定理或尖锐常数证明。
% 下文势表示、有限指数卷积估计与级数求和均在本书展开；
% 只引用实读首页作历史定位与进一步阅读，不声称复现Moser最佳常数。

\subsection{\texorpdfstring{\(p=d\)}{p=d}}
\label{b1:subsec:240501}

为了在卷积估计中让输出指数增大，使用\cref{b1:lem:young-convolution-finite}的有限指数 Young 卷积不等式。该基础工具已在卷积理论建立处证明，不把它的唯一正式出处留在本选读节中。

\begin{proposition}\label{b1:prop:critical-finite-q-estimate}
设 \(\Omega\subset\mathbb R^d\) 为非空有界开集，\(D=\operatorname{diam}\Omega\)，且 \(d\ge2\)。对 \(d\le q<\infty\)、\(u\in W_0^{1,d}(\Omega)\)，有
\begin{equation}\label{b1:eq:critical-q-growth}
 \|u\|_{L^q(\Omega)}
 \le C_dD^{d/q}q^{1-1/d}\|\nabla u\|_{L^d(\Omega)}.
\end{equation}
其中 \(C_d\) 不依赖 \(q,u,\Omega\)。
\end{proposition}
\begin{proof}
先取 \(u\in C_c^\infty(\Omega)\)，在全空间补零。固定 \(x\) 及单位向量 \(\omega\)，沿射线积分得
\[
 u(x)=-\int_0^\infty\nabla u(x+r\omega)\cdot\omega\,\mathrm dr.
\]
对单位球面平均，再使用极坐标换元，得到
\[
 |u(x)|\le\frac1{d\omega_d}
        \int_{\mathbb R^d}\frac{|\nabla u(y)|}{|x-y|^{d-1}}\,\mathrm dy.
\]
若 \(x\in\Omega\)，只有 \(y\in\Omega\) 对积分有贡献，故 \(|x-y|\le D\)。记
\[
 K_D(h)=|h|^{1-d}\mathbf1_{\{|h|\le D\}},\qquad
 b_q=\frac{dq}{(d-1)q+d}.
\]
则 \(1\le b_q<d/(d-1)\)，并且
\[
 1+\frac1q=\frac1d+\frac1{b_q}.
\]
Young 卷积不等式给出
\[
 \|u\|_{L^q(\Omega)}
 \le C_d\|K_D\|_{L^{b_q}}\cdot\|\nabla u\|_{L^d(\Omega)}.
\]
核的范数可以直接计算：
\[
 \begin{aligned}
 \|K_D\|_{L^{b_q}}
 &=\left(\frac{d\omega_d}
                    {d-(d-1)b_q}\right)^{1/b_q}
        D^{\,d/b_q-(d-1)}\\
 &\le C_d q^{1/b_q}D^{d/q}
 \le C_d q^{1-1/d}D^{d/q}.
 \end{aligned}
\]
这里 \(d/b_q-(d-1)=d/q\)，
\(d-(d-1)b_q=db_q/q\)，且
\(q^{1/b_q}=q^{1-1/d}q^{1/q}\)，最后一因子在 \(q\ge1\) 上有界。这解释了常数中 \(q^{1-1/d}\) 的来源。

对一般 \(u\in W_0^{1,d}(\Omega)\)，取内部紧支集光滑逼近。将已证估计用于逼近列两项之差，可知它在每个固定有限 \(L^q\) 中为 Cauchy 列；其极限与原来的 \(L^d\) 极限一致，所以就是 \(u\)。取极限完成证明。
\end{proof}

在有界 Lipschitz 区域上，一般 \(u\in W^{1,d}(\Omega)\) 也属于每个有限 \(L^q(\Omega)\)。事实上，\cref{b1:thm:lipschitz-domain-extension}给出支集留在固定紧集中的全空间延拓；把这个紧集放进一个更大的球，延拓函数属于该球的 \(W_0^{1,d}\)，便可应用上面的估计。所得界包含 \(\|u\|_{W^{1,d}(\Omega)}\)，不能只用梯度，因为非零常数也在此空间中。

\ProseParagraphBreak

不能在\eqref{b1:eq:critical-q-growth}中令 \(q=\infty\)：常数随 \(q\) 增大而发散。接下来的求和恰好利用这一增长速度，而不是把它忽略。

\subsection{Trudinger–Moser 不等式}
\label{b1:subsec:240502}

记 \(d'=d/(d-1)\)。\term[trudinger-moser-budengshi]{Trudinger--Moser 不等式}[Trudinger--Moser inequality]的一个非最优常数版本如下。

\begin{theorem}[Trudinger--Moser，非最优常数版]\label{b1:thm:trudinger-moser}
设 \(d\ge2\)，\(\Omega\subset\mathbb R^d\) 为非空有界开集，\(D=\operatorname{diam}\Omega\)。存在只依赖于 \(d\) 的正数 \(c_d,C_d\)，使每个 \(u\in W_0^{1,d}(\Omega)\) 满足
\begin{equation}\label{b1:eq:trudinger-moser}
 \int_\Omega
  \exp\left(c_d
       \left|\frac{u(x)}{\|\nabla u\|_{L^d(\Omega)}}\right|^{d'}\right)
       \mathrm dx
 \le C_dD^d
\end{equation}
当 \(\|\nabla u\|_d>0\)。若梯度范数为零，则 \(u=0\)，将被积函数解释为 \(1\)。这里不声称 \(c_d\) 是最大的允许常数。
\end{theorem}
\begin{proof}
由零边界 Poincaré 不等式，梯度为零时 \(u=0\)，所列约定成立。其余情形可通过数乘归一化为 \(\|\nabla u\|_d=1\)。展开非负指数级数，单调收敛定理给出
\[
 \int_\Omega e^{c|u|^{d'}}
 =|\Omega|+
   \sum_{j=1}^\infty\frac{c^j}{j!}
          \|u\|_{L^{d'j}(\Omega)}^{d'j}.
\]
对 \(d'j\ge d\)，应用\eqref{b1:eq:critical-q-growth}，
\[
 \|u\|_{L^{d'j}}^{d'j}
 \le D^d C_d^{d'j}(d'j)^j.
\]
由
\[
 \log(j!)=\sum_{\ell=1}^j\log\ell
 \ge\int_1^j\log t\,\mathrm dt
 =j\log j-j+1
\]
得到 \(j!\ge(j/e)^j\)。因此上述级数的相应项不超过
\[
 D^d\bigl(c\,C_d^{d'}d'e\bigr)^j.
\]
选取只依赖于 \(d\) 的充分小 \(c=c_d>0\)，使括号不超过 \(1/2\)，这些高次项之和便受常数倍的 \(D^d\) 控制。

剩下满足 \(d'j<d\) 的项只有有限个。由 Hölder 不等式及零边界 Poincaré 估计，
\[
 \|u\|_{L^{d'j}}^{d'j}
 \le|\Omega|^{\,1-d'j/d}\|u\|_{L^d}^{d'j}
 \le C_dD^d.
\]
这里把 \(\Omega\) 放入直径不超过常数倍 \(D\) 的球，使用
\(|\Omega|\le\omega_dD^d\) 与 \(\|u\|_d\le C_dD\)。
有限低次项、常数项和高次几何级数合起来，给出所需估计。最后恢复原函数的梯度范数，完成证明。
\end{proof}

这个定理是对梯度范数有界函数族的统一估计，不只是说每个固定函数在某个很小参数下可积。指数幂 \(d'\) 使高次 \(L^q\) 常数产生的 \(j^j\) 与阶乘的增长相抵；若把指数幂再增大，后者便未必能够控制前者。习题中的集中函数列将直接说明这种更强统一估计为何失败。

\ProseParagraphBreak

一般 \(W^{1,d}\) 函数若没有零边界条件，不能在分母中只放梯度；常数函数已足以显示问题。可在有界连通 Lipschitz 区域上先减去均值，再使用 Poincaré 和固定支集延拓取得相应的指数估计，但其常数会随区域几何改变。本文所证明的\eqref{b1:eq:trudinger-moser}保留了一个便于检查的版本：任意有界开集、零边界闭包、非最优常数和明确的直径尺度。

\subsection{Morrey–Campanato 观点}
\label{b1:subsec:240503}

Morrey 证明中的决定性步骤不是逐点微分，而是球平均的振幅随半径衰减。把导数暂时拿掉，只保留这种平均条件，就得到 Campanato 观点的基本形式。

\begin{proposition}[Campanato 型平均条件]\label{b1:prop:campanato-holder}
设 \(u\in L^1_{\mathrm{loc}}(\mathbb R^d)\)、\(0<\alpha\le1\)。若存在 \(A<\infty\)，使每个球 \(B(x,r)\) 都满足
\begin{equation}\label{b1:eq:campanato-mean-condition}
 \frac1{|B(x,r)|}\int_{B(x,r)}
       |u-u_{B(x,r)}|\,\mathrm dy\le Ar^\alpha,
\end{equation}
则 \(u\) 有唯一连续代表 \(u^*\)，并有
\[
 |u^*(x)-u^*(y)|\le C_{d,\alpha}A|x-y|^\alpha.
\]
反过来，任意满足全域 \(\alpha\)-Hölder 估计的函数都满足相应的平均条件。
\end{proposition}
\begin{proof}
固定 \(x,R\)。两相邻同心球的体积比为 \(2^d\)，故
\[
 |u_{B(x,2^{-j-1}R)}-u_{B(x,2^{-j}R)}|
 \le2^dA(2^{-j}R)^\alpha.
\]
右侧为可和几何级数，所以球均值在每个点有有限极限，并有
\[
 |u^*(x)-u_{B(x,R)}|
 \le C_{d,\alpha}AR^\alpha.
\]
任意半径可夹在两个相邻二分半径之间，使用相同的体积比估计，说明极限不依赖这列半径。Lebesgue 微分定理保证 \(u^*=u\) 几乎处处。

对 \(x\ne y\)，令 \(R=|x-y|\)。两个球 \(B(x,R)\)、\(B(y,R)\) 都包含在 \(B(x,3R)\) 内，且体积比受维数控制。因此各自均值与大球均值之差均不超过 \(C_dAR^\alpha\)。再加上两个点到各自均值的误差，得到所求双点估计。它同时证明连续性；两个连续代表若几乎处处相等，就处处相等，因此代表唯一。

反过来，若 \(|u(z)-u(w)|\le L|z-w|^\alpha\)，则
\[
 \begin{aligned}
 \frac1{|B|}\int_B|u-u_B|
 &\le\frac1{|B|^2}\int_B\int_B|u(z)-u(w)|\,\mathrm dz\,\mathrm dw\\
 &\le L(2r)^\alpha,
 \end{aligned}
\]
这就是平均条件。
\end{proof}

更常见的表述把左侧换成 \(q\) 次平均振幅，\(1\le q<\infty\)。它由 Hölder 不等式控制这里的一次平均，而 Hölder 连续函数也控制每个这样的 \(q\) 次平均，因此在正衰减指数范围内仍刻画相同的连续性。一般 Campanato 空间允许更多参数和多项式近似；这里仅需平均常数的基本情形。

\ProseParagraphBreak

在临界指数 \(p=d\) 时，Poincaré 不等式只给出
\[
 \frac1{|B|}\int_B|u-u_B|
 \le C_d\left(\int_B|\nabla u|^d\right)^{1/d}.
\]
右侧没有固定的正半径幂。对固定的可积梯度，它会随球的体积缩小而趋零，但这些误差未必沿二分尺度可求和。因而“平均振幅趋零”与“可以把球均值求和为连续代表”仍有距离；仅把前一命题的 \(\alpha\) 设为零，会破坏其关键级数。

\subsection{临界现象}
\label{b1:subsec:240504}

下面构造一个 \(W_0^{1,d}\) 中本性无界的函数，说明临界指数的确不能一般给出连续或有界代表。取
\[
 0<\beta<1-\frac1d,
\]
并选光滑径向截断 \(\chi\)，使它在原点附近等于 \(1\)、支集包含于 \(B(0,1/2)\)。对 \(0<|x|<1/2\)，令
\[
 u(x)=\chi(x)\bigl(\log(e/|x|)\bigr)^\beta,
\]
在支集外取零，在原点任意赋一个有限值。显然它在每个原点邻域中本性无界。

\ProseParagraphBreak

函数本身属于 \(L^d(B(0,1))\)，因为径向体积因子 \(r^{d-1}\) 控制任意固定的对数幂。在 \(\chi=1\) 的小球内，梯度长度为
\[
 |\nabla u(x)|
 =\frac{\beta}{|x|}
       \bigl(\log(e/|x|)\bigr)^{\beta-1},
\]
故其 \(d\) 次积分等价于
\[
 \int_0^\varepsilon
      \frac{\mathrm dr}
           {r\bigl(\log(e/r)\bigr)^{d(1-\beta)}}.
\]
代换 \(t=\log(e/r)\) 后，这是尾部的
\(\int^\infty t^{-d(1-\beta)}\,\mathrm dt\)，由于 \(d(1-\beta)>1\) 而收敛。截断的变化只发生在避开原点的紧环带上，不增加奇性。

\ProseParagraphBreak

还须确认这些经典导数确实是弱导数。对测试函数在挖去半径 \(\varepsilon\) 小球的区域上分部积分，内球面边界项的绝对值受
\[
 C\varepsilon^{d-1}
       \bigl(\log(e/\varepsilon)\bigr)^\beta
       \|\varphi\|_\infty
\]
控制，因 \(d\ge2\) 而趋零。函数和候选导数局部可积，故令 \(\varepsilon\downarrow0\) 得到全域弱导数恒等式。于是 \(u\in W^{1,d}\)，又因为它具有内部紧支集，有限指数的内部逼近给出 \(u\in W_0^{1,d}(B(0,1))\)。

\ProseParagraphBreak

它没有有界代表，但仍服从本节证明的指数估计。这并不矛盾：指数可积性限制的是大值集合的体积，而非禁止函数取得任意大的值。另一方面，当 \(d=1\) 时，\(W^{1,1}\) 的绝对连续代表已经由积分基本定理得到有界控制，本节的临界反例和指数幂 \(d/(d-1)\) 都不适用。

\ProseParagraphBreak

本章至此建立了梯度对均值振幅、可积性与连续性的不同控制。下一章将询问一个新的问题：若一列函数都服从这些控制，能否从中抽出强收敛子列？单个函数的嵌入估计不会自动给出序列紧性；还须排除质量逃向无穷远、平移失控和临界尺度的集中。


\begin{chaptersummary}
梯度首先只能控制函数相对于常数的变化。均值与零边界条件以不同方式消除了这一自由度：前者在一般区域上需要连通性及合适的几何控制，后者可以在任意有界开集上通过紧支集逼近建立。

指数低于维数时，逐坐标积分与幂次代换提高了函数值的可积性；指数高于维数时，可和的球平均误差产生了连续代表。两条路线都保留了尺度，并各有不能直接纳入的端点。高阶与分数阶嵌入需要把这些工具实际作用于相应导数或差分，不能仅对定理中的阶数作形式替换。

临界指数没有一般的 \(L^\infty\) 控制，却仍能给出指数可积性。常数的增长、平均振幅的衰减和集中反例共同说明了这种边界。下一章将把注意力从单个函数转到函数列，研究这些有界控制在什么条件下能够产生强收敛子列。
\end{chaptersummary}

\BookExercises

% \chapref{b1:ch:24}习题临时稿；不另设 BookExercises 入口。
% 八题均为自拟的正文工具变式；各解答展开新增计算或构造，不冒认教材原题编号。
% 依赖止于本章，不调用\chapref{b1:ch:25}的紧嵌入或平移紧性判据。

\begin{exercise}\label{b1:ex:24-componentwise-poincare}
设 \(1\le p<\infty\)，\(\Omega_1,\ldots,\Omega_N\subset\mathbb R^d\) 为非空、有界、连通的 Lipschitz 开集，且它们的闭包两两不交。令
\(\Omega=\bigcup_{j=1}^N\Omega_j\)，并对 \(u\in W^{1,p}(\Omega)\) 定义
\[
 Pu=\sum_{j=1}^Nu_{\Omega_j}\mathbf1_{\Omega_j},
 \quad u_{\Omega_j}=\frac1{|\Omega_j|}\int_{\Omega_j}u.
\]
\begin{enumerate}
 \item\label{b1:item:24-componentwise-estimate}
 证明 \(P\) 是 \(L^p(\Omega)\) 上的压缩投影，且
 \[
  \|u-Pu\|_{L^p(\Omega)}
       \le C\|\nabla u\|_{L^p(\Omega)}.
 \]
 若各分量均值 Poincaré 不等式的最优常数为 \(C_j\)，证明此处最优常数为
 \(\max_j C_j\)。
 \item\label{b1:item:24-componentwise-kernel}
 刻画 \(W^{1,p}(\Omega)\) 中满足 \(\nabla u=0\) 的全部函数，并证明当 \(N\ge2\) 时，不能把上式中的 \(Pu\) 换成全域平均值 \(u_\Omega\)。
\end{enumerate}
\end{exercise}
% 来源：自拟；有限分量的直和估计及梯度核，补充连通域均值 Poincaré 的假设边界。
\begin{solution}
由 Jensen 不等式，
\[
 |\Omega_j|\cdot|u_{\Omega_j}|^p
       \le\int_{\Omega_j}|u|^p.
\]
求和得到 \(\|Pu\|_p\le\|u\|_p\)。取过一次分量平均后，每个分量上已是常函数，所以 \(P^2=P\)；线性也由积分的线性得到。

分量上的常函数在 \(\Omega\) 内没有弱导数。事实上，测试函数的支集与
\(\partial\Omega\) 保持正距离，在每个分量内积分其偏导都为零；这里没有跨越域外边界的跳跃项。因此 \(Pu\in W^{1,p}(\Omega)\)，且
\(\nabla(u-Pu)=\nabla u\)。在每个分量应用定理%
\ref{b1:thm:connected-domain-poincare}，再把 \(p\) 次方相加，得到
\[
 \begin{aligned}
 \|u-Pu\|_p^p
 &=\sum_{j=1}^N\|u-u_{\Omega_j}\|_{L^p(\Omega_j)}^p\\
 &\le\left(\max_j C_j\right)^p
                       \sum_{j=1}^N\|\nabla u\|_{L^p(\Omega_j)}^p.
 \end{aligned}
\]
最优常数的下界可以逐个分量测试：在 \(\Omega_j\) 上任选一个函数，在其他分量上置零，整体比值恰好成为该分量的比值。取上确界并遍历 \(j\)，就得到所求的精确常数，不要求最优值由某个函数取得。

若 \(\nabla u=0\)，刚才的估计给出 \(u=Pu\) 几乎处处，故 \(u\) 在每个分量上分别为常数；反向已在前面说明。若 \(N\ge2\)，取
\[
 u=|\Omega_2|\mathbf1_{\Omega_1}
           -|\Omega_1|\mathbf1_{\Omega_2},
\]
并在其余分量上置零。此函数的全域平均为零、梯度为零，却不是零函数。因此全域平均不能消除梯度核中的全部自由度；需要消除的是每个分量上的一个常数。
\end{solution}

\begin{exercise}\label{b1:ex:24-thin-neck-constant}
设 \(d\ge2\)、\(1\le p<\infty\)、\(0<\varepsilon<1/2\)。考虑两个盒子由细颈连接而成的开集
\[
 \begin{aligned}
 \Omega_\varepsilon={}&
   \bigl((-2,-1/2)\times(-1,1)^{d-1}\bigr)\\
 &{}\cup\bigl((-1,1)\times(-\varepsilon,\varepsilon)^{d-1}\bigr)\\
 &{}\cup\bigl((1/2,2)\times(-1,1)^{d-1}\bigr).
 \end{aligned}
\]
令 \(C_\varepsilon\) 为该连通 Lipschitz 开集上均值 Poincaré 不等式的最优常数。证明
\[
 C_\varepsilon\ge
       \frac{3^{1/p}}2\,\varepsilon^{-(d-1)/p}.
\]
据此说明：对全部有界连通 Lipschitz 开集，常数不能只由维数、指数以及直径和体积的统一上下界控制。
\end{exercise}
% 来源：自拟；固定体积盒子与退化细颈的显式测试，不借用特征值或谱理论。
\begin{solution}
定义分段线性函数
\[
 h(s)=
 \begin{cases}
  -1,&s\le-1/2,\\
  2s,&-1/2<s<1/2,\\
  1,&s\ge1/2,
 \end{cases}
 \quad u_\varepsilon(x)=h(x_1).
\]
这个函数在全空间 Lipschitz，因而其在 \(\Omega_\varepsilon\) 上的限制属于
\(W^{1,p}\)。区域关于 \(x_1\mapsto-x_1\) 对称，而函数关于此变换变号，所以
\((u_\varepsilon)_{\Omega_\varepsilon}=0\)。

两个大盒子上的函数模都等于 \(1\)。它们各自的体积为
\((3/2)2^{d-1}\)，故
\[
 \|u_\varepsilon\|_{L^p(\Omega_\varepsilon)}^p
              \ge3\cdot2^{d-1}.
\]
唯一非零的偏导是中央细颈上的 \(\partial_1u_\varepsilon=2\)。这一部分的长度为 \(1\)，横截面体积为 \((2\varepsilon)^{d-1}\)，因而
\[
 \|\nabla u_\varepsilon\|_{L^p(\Omega_\varepsilon)}^p
              =2^p(2\varepsilon)^{d-1}.
\]
将这个函数代入最优常数的定义，即得
\[
 C_\varepsilon^p\ge
     \frac{3\cdot2^{d-1}}{2^p(2\varepsilon)^{d-1}}
     =\frac3{2^p}\varepsilon^{-(d-1)}.
\]

这些区域均包含固定的两个大盒子，又都包含在
\((-2,2)\times(-1,1)^{d-1}\) 中，故体积和直径均有与 \(\varepsilon\) 无关的正下界及有限上界。对每个固定 \(\varepsilon\)，盒子的平面边界与有限个转角给出 Lipschitz 边界；开集连通，因为细颈与两个盒子分别重叠。随 \(\varepsilon\) 缩小，退化的是连接几何，而不是这些粗略尺寸。上面的下界趋于无穷，证明了最后的断言。
\end{solution}

\begin{exercise}\label{b1:ex:24-linear-normalization}
设 \(\Omega\subset\mathbb R^d\) 为非空、有界、连通的 Lipschitz 开集，
\(1\le p<\infty\)。本题取实值函数，并令
\(\ell:W^{1,p}(\Omega)\to\mathbb R\) 为连续线性泛函。证明下列条件等价。
\begin{equivalences}
 \item\label{b1:item:24-normalization-detects-constants}
 \(\ell(1)\ne0\)。
 \item\label{b1:item:24-normalization-poincare}
 存在有限常数 \(C_\ell\)，使所有满足 \(\ell(u)=0\) 的 \(u\) 都有
 \(\|u\|_{L^p(\Omega)}\le C_\ell\|\nabla u\|_{L^p(\Omega)}\)。
\end{equivalences}
进一步，若 \(A\subset\Omega\) 可测且 \(|A|>0\)，证明条件 \(\int_Au=0\) 蕴含
\[
 \|u\|_{L^p(\Omega)}
 \le C_P\Bigl(1+(|\Omega|/|A|)^{1/p}\Bigr)
          \|\nabla u\|_{L^p(\Omega)},
\]
其中 \(C_P\) 是通常均值 Poincaré 不等式的任一个可用常数。
\end{exercise}
% 来源：自拟；用一般连续线性约束代替全域均值，再给出可测子集上的明确常数。
\begin{solution}
若 \(\ell(1)=0\)，常函数 \(u=1\) 就属于 \(\ker\ell\)，但其梯度为零、函数范数非零，所以第二个条件不可能成立。

现在设 \(\ell(1)\ne0\)，并取 \(\ell(u)=0\)。写成
\[
 u=v+c,\quad c=u_\Omega,\quad v=u-u_\Omega.
\]
由\cref{b1:thm:connected-domain-poincare}及弱导数范数的有限维等价性，有
\[
 \|v\|_{L^p(\Omega)}\le C_P\|\nabla u\|_p,
 \quad \|v\|_{W^{1,p}(\Omega)}\le C_0\|\nabla u\|_p,
\]
其中 \(C_0\) 只依赖于 \(\Omega,d,p\)。因为
\(0=\ell(v)+c\ell(1)\)，连续性给出
\[
 |c|\le\frac{\|\ell\|}{|\ell(1)|}
                   \|v\|_{W^{1,p}(\Omega)}
       \le\frac{C_0\|\ell\|}{|\ell(1)|}\|\nabla u\|_p.
\]
于是
\[
 \|u\|_p\le\|v\|_p+|\Omega|^{1/p}\cdot|c|
 \le\left(C_P+
       \frac{|\Omega|^{1/p}C_0\|\ell\|}{|\ell(1)|}\right)
                 \|\nabla u\|_p,
\]
这就证明等价性。

对于最后的条件，仍用同一分解；从 \(\int_Au=0\) 得到
\[
 |c|=\left|\frac1{|A|}\int_Av\right|
       \le |A|^{-1/p}\|v\|_{L^p(A)}.
\]
这里 \(p=1\) 时直接使用积分三角不等式，\(p>1\) 时使用 Hölder 不等式。因此
\[
 \|u\|_p\le
     \Bigl(1+(|\Omega|/|A|)^{1/p}\Bigr)
                   \|v\|_p,
\]
再代入通常的 Poincaré 估计，便得到所列常数。这个约束不要求 \(A\) 是开集，也不要求它包含一个球。
\end{solution}

\begin{exercise}\label{b1:ex:24-supercritical-algebra}
设 \(\Omega\subset\mathbb R^d\) 为非空有界 Lipschitz 开集，\(d<p<\infty\)，本题取实值函数。
\begin{enumerate}
 \item\label{b1:item:24-sobolev-algebra}
 证明乘法满足
 \[
  \|uv\|_{W^{1,p}(\Omega)}
       \le C\|u\|_{W^{1,p}(\Omega)}
                 \cdot\|v\|_{W^{1,p}(\Omega)}.
 \]
 由此证明，适当放大原范数后，新范数仍完备且满足次乘性
 \(\|uv\|_{\mathrm A}\le\|u\|_{\mathrm A}\cdot\|v\|_{\mathrm A}\)，常函数 \(1\) 为乘法单位元。
 \item\label{b1:item:24-sobolev-inversion}
 若 \(u\in W^{1,p}(\Omega)\) 且 \(|u|\ge\delta>0\) 几乎处处，证明
 \(1/u\in W^{1,p}(\Omega)\)，并给出它的函数项与梯度项估计。进一步证明：在这类远离零的函数中，若 \(u_n\to u\) 于 \(W^{1,p}\)，且所有 \(|u_n|\ge\delta\)，则
 \(1/u_n\to1/u\) 于 \(W^{1,p}\)。
\end{enumerate}
\end{exercise}
% 来源：自拟；Morrey 嵌入与已建立的 Leibniz、链式法则组合，强调代数性所需的上确界控制。
\begin{solution}
由\cref{b1:cor:morrey-domain}，
\[
 \|w\|_\infty\le C_M\|w\|_{W^{1,p}(\Omega)}.
\]
所以 \(u,v\) 都本性有界，可以使用\cref{b1:prop:sobolev-product}的乘积公式。
函数项和梯度项分别满足
\[
 \|uv\|_p\le\|u\|_\infty\cdot\|v\|_p,
 \quad
 \|\nabla(uv)\|_p
 \le\|u\|_\infty\cdot\|\nabla v\|_p
       +\|v\|_\infty\cdot\|\nabla u\|_p.
\]
代入上确界估计并合并各项，就得到第一问。取一个不小于 \(1\) 的可用常数 \(C\)，将范数改为
\(\|w\|_{\mathrm A}=C\|w\|_{W^{1,p}}\)，则
\(\|uv\|_{\mathrm A}\le\|u\|_{\mathrm A}\cdot\|v\|_{\mathrm A}\)。
等价范数保持完备性，而有界域上的常函数 \(1\) 属于这个空间，并且是乘法单位元。

对第二问，取光滑函数 \(\zeta\)，使它在 \([-\delta/2,\delta/2]\) 上为零，在
\((-\infty,-\delta]\cup[\delta,\infty)\) 上为一，令
\(F(t)=\zeta(t)/t\) 并规定 \(F(0)=0\)。这样 \(F\) 在全实轴光滑，导数有界，且
\(F(u)=1/u\)。由\cref{b1:thm:sobolev-chain}，
\[
 \nabla(1/u)=-u^{-2}\nabla u
       \quad\text{几乎处处},
\]
且
\[
 \|1/u\|_p\le|\Omega|^{1/p}\delta^{-1},
 \quad
 \|\nabla(1/u)\|_p\le\delta^{-2}\|\nabla u\|_p.
\]
这里有界域使函数项可积，不仅仅是导数项可积。

最后，Morrey 估计使 \(u_n\to u\) 也成为连续代表的一致收敛。倒数之差满足
\[
 \|1/u_n-1/u\|_p\le\delta^{-2}\|u_n-u\|_p.
\]
在集合 \(\{\,t:|t|\ge\delta\,\}\) 上，映射 \(t\mapsto t^{-2}\) 满足
\[
 |a^{-2}-b^{-2}|\le2\delta^{-3}|a-b|.
\]
当 \(a,b\) 同号时这来自导数界；异号时可先比较 \(|a|,|b|\)，并使用
\(\bigl||a|-|b|\bigr|\le|a-b|\)。因此
\[
 \begin{aligned}
 \|\nabla(1/u_n)-\nabla(1/u)\|_p
 &\le\delta^{-2}\|\nabla u_n-\nabla u\|_p\\
 &\quad+2\delta^{-3}\|u_n-u\|_\infty
                         \cdot\|\nabla u\|_p
 \longrightarrow0.
 \end{aligned}
\]
两项合起来便给出所求 Sobolev 范数收敛。
\end{solution}

\begin{exercise}\label{b1:ex:24-morrey-logarithmic-cusp}
设 \(d\ge2\)、\(d<p<\infty\)，并记 \(\alpha=1-d/p\)。
对于实数 \(\gamma\)，在单位球 \(B\) 上定义
\[
 u(0)=0,\quad
 u(x)=|x|^\alpha\bigl(\log(e/|x|)\bigr)^{-\gamma}
                  \quad(0<|x|<1).
\]
证明
\[
 u\in W^{1,p}(B)\quad\Longleftrightarrow\quad\gamma p>1.
\]
在 \(\gamma p>1\) 时，证明 \(u\) 的连续代表属于 \(C^{0,\alpha}(\overline B)\)，但不属于任何
\(C^{0,\beta}(\overline B)\)，其中 \(\alpha<\beta\le1\)。
这里要检验一个固定函数的正则性，不能只用函数族的缩放排除统一估计。
\end{exercise}
% 来源：自拟；临界幂次乘对数因子的可积性计算，与正文的统一 Hölder 估计尖锐性区分。
\begin{solution}
令 \(L(r)=\log(e/r)\)。由于任意正幂的衰减快于对数的增长，
\(r^\alpha L(r)^{-\gamma}\to0\)，所以给定函数在 \(\overline B\) 连续，特别属于
\(L^p(B)\)。在原点以外，径向求导得到
\[
 \frac{\mathrm d}{\mathrm dr}
       \bigl(r^\alpha L(r)^{-\gamma}\bigr)
 =r^{\alpha-1}L(r)^{-\gamma}
                     \left(\alpha+\frac\gamma{L(r)}\right).
\]
当 \(r\) 足够小时，最后括号中的量处于两个严格为正的常数之间。因此经典梯度的
\(p\) 次积分在原点附近有限，当且仅当
\[
 \int_0^{1/2}r^{d-1+p(\alpha-1)}L(r)^{-\gamma p}\,\mathrm dr
 =\int_0^{1/2}\frac{\mathrm dr}{rL(r)^{\gamma p}}<\infty.
\]
代换 \(t=L(r)\) 后，这等价于 \(\int_{\log(2e)}^\infty t^{-\gamma p}\,\mathrm dt<\infty\)，也就是 \(\gamma p>1\)。
如果 \(u\in W^{1,p}(B)\)，它的弱梯度在不含原点的区域必与这个经典梯度一致，故该条件必要。

为证明充分性，还须核实原点没有遗漏的分布项。设 \(\gamma p>1\)，取光滑
\(\eta:[0,\infty)\to[0,1]\)，使它在 \([0,1]\) 上为零，在 \([2,\infty)\) 上为一，令
\[
 u_\varepsilon(x)=\eta(|x|/\varepsilon)u(x).
\]
每个 \(u_\varepsilon\) 在 \(B\) 内光滑。记上面算出的经典梯度为 \(g\)，并在原点随意赋值。已经知道 \(g\in L^p(B)\)，且
\[
 \nabla u_\varepsilon
  =\eta(|x|/\varepsilon)g
       +\varepsilon^{-1}\eta'(|x|/\varepsilon)
                         u(x)\frac{x}{|x|}
\]
在原点以外成立。第一项在 \(L^p\) 中趋于 \(g\)。第二项仅支于
\(\{\varepsilon<|x|<2\varepsilon\}\)，其范数的 \(p\) 次方不超过
\[
 C\varepsilon^{-p}\varepsilon^{d+\alpha p}
                  L(\varepsilon)^{-\gamma p}
   =C L(\varepsilon)^{-\gamma p}\longrightarrow0.
\]
这里用到了环带上 \(L(|x|)\) 与 \(L(\varepsilon)\) 的可比性，以及
\(d+\alpha p-p=0\)。同时 \(u_\varepsilon\to u\) 于 \(L^p\)。在弱导数恒等式中通过极限，得到 \(u\in W^{1,p}(B)\) 且 \(\nabla u=g\)。

由\cref{b1:cor:morrey-domain}，这个空间的元素有 \(\alpha\)-Hölder 连续代表。它与题中已有的连续函数几乎处处相同，因而处处相同。可是对任意 \(\beta>\alpha\)，
\[
 \frac{|u(re_1)-u(0)|}{|re_1|^\beta}
       =r^{\alpha-\beta}L(r)^{-\gamma}
                  \longrightarrow\infty.
\]
所以任何更大的 Hölder 指数都不能用于这个固定函数。对数衰减使临界梯度可积，却不能补偿正幂之间的差距。
\end{solution}

\begin{exercise}\label{b1:ex:24-moser-concentration}
设 \(d\ge2\)，记 \(\omega_d=|B(0,1)|\)、\(\sigma_{d-1}=d\omega_d\)，固定 \(R=1/2\)。对 \(L\ge1\)，在 \(\mathbb R^d\) 上定义
\[
 u_L(x)=\sigma_{d-1}^{-1/d}
 \begin{cases}
  L^{(d-1)/d},& |x|\le R e^{-L},\\
  L^{-1/d}\log(R/|x|),& R e^{-L}<|x|<R,\\
  0,&|x|\ge R.
 \end{cases}
\]
\begin{enumerate}
 \item\label{b1:item:24-moser-norms}
 证明 \(u_L\in W_0^{1,d}(B(0,1))\)，并且
 \(\|\nabla u_L\|_d=1\)、\(\|u_L\|_d\to0\)。
 \item\label{b1:item:24-moser-supercritical-exponent}
 证明若 \(q>d/(d-1)\)，则对于每个 \(a>0\)，都有
 \[
  \int_{B(0,1)}\exp(a|u_L|^q)\,\mathrm dx
                       \longrightarrow\infty.
 \]
 在 \(q=d/(d-1)\) 时，进一步证明同样的发散对于
 \(a>d\,\sigma_{d-1}^{1/(d-1)}\) 成立。
\end{enumerate}
由此说明超临界指数不可能具有梯度单位球上的统一指数积分界；本题不要求证明临界系数上界能够取得。
\end{exercise}
% 来源：自拟整理经典径向对数集中机制；这里逐项证明所需结论，不冒认未实读的 Moser 原题或最优常数证明。
\begin{solution}
对于每个固定 \(L\)，径向剖面连续且分段光滑，导数有界。因此 \(u_L\) 是全空间的
Lipschitz 函数，支集包含于 \(\overline{B(0,R)}\Subset B(0,1)\)。它属于
\(W^{1,d}\)，由\cref{b1:prop:compact-sobolev-smooth-approximation}的内部紧支集逼近，得到 \(u_L\in W_0^{1,d}(B(0,1))\)。

梯度只在中间环带非零，其模为
\(\sigma_{d-1}^{-1/d}L^{-1/d}|x|^{-1}\)。极坐标积分给出
\[
 \int_{\mathbb R^d}|\nabla u_L|^d\,\mathrm dx
       =L^{-1}\int_{Re^{-L}}^R\frac{\mathrm dr}{r}=1.
\]
对函数本身，在环带使用 \(t=\log(R/r)\) 代换，得到
\[
 \|u_L\|_d^d
 =\frac{\omega_dR^d}{\sigma_{d-1}}L^{d-1}e^{-dL}
       +\frac{R^d}{L}\int_0^L t^d e^{-dt}\,\mathrm dt.
\]
第一个项趋于零；第二个项也趋于零，因为
\(\int_0^\infty t^d e^{-dt}\,\mathrm dt<\infty\)。这证明第一问的两个范数结论。

内球 \(B(0,Re^{-L})\) 上的函数值是常数，故
\[
 \int_{B(0,1)}\exp(a|u_L|^q)\,\mathrm dx
 \ge\omega_dR^d
       \exp\left(-dL+
          a\sigma_{d-1}^{-q/d}L^{q(d-1)/d}\right).
\]
若 \(q>d/(d-1)\)，最后的正项关于 \(L\) 的幂次严格大于 \(1\)，因而它超过线性负项，右端趋于无穷。若 \(q=d/(d-1)\)，指数变为
\[
 \left(a\sigma_{d-1}^{-1/(d-1)}-d\right)L.
\]
题设对 \(a\) 的条件恰使括号为正，所以仍然发散。这给出了临界指数下系数的一个必要上界；证明没有给出该界处的统一可积估计。
\end{solution}

\begin{exercise}\label{b1:ex:24-higher-order-scaling}
设 \(B=B(0,1)\subset\mathbb R^d\)、\(k\ge2\) 为整数、\(1\le p<\infty\)。本题只测试
\(C_c^\infty(B)\) 中的函数，不使用高阶延拓定理。
\begin{enumerate}
 \item\label{b1:item:24-higher-order-lq-obstruction}
 当 \(kp<d\) 时，设 \(1\le q\le\infty\)。证明若
 \[
  \|u\|_{L^q(B)}\le C\|u\|_{W^{k,p}(B)}
       \quad\text{对全部 }u\in C_c^\infty(B)
 \]
 成立，则必须有 \(q\le dp/(d-kp)\)。这里允许 \(q=\infty\)。
 \item\label{b1:item:24-higher-order-holder-obstruction}
 设 \(0\le j<k\)、\(0<\beta\le1\)。若存在有限常数 \(C\)，使所有
 \(u\in C_c^\infty(B)\) 都满足
 \[
  \max_{|\kappa|=j}[D^\kappa u]_{0,\beta;\overline B}
        \le C\|u\|_{W^{k,p}(B)},
 \]
 证明必要关系为 \(j+\beta\le k-d/p\)。
\end{enumerate}
说明为什么第一问的同一种集中计算，在 \(kp=d\) 时不能单独排除目标空间
\(L^\infty\)。
\end{exercise}
% 来源：自拟；高阶非齐次范数中各阶项的尺度分离，以及高阶 Hölder 目标的必要限制。
\begin{solution}
取非零 \(\varphi\in C_c^\infty(B)\)，并对 \(0<\varepsilon<1\) 令
\[
 u_\varepsilon(x)=\varepsilon^{k-d/p}\varphi(x/\varepsilon).
\]
它的支集包含于 \(\varepsilon\operatorname{supp}\varphi\Subset B\)。对每个
\(|\kappa|\le k\)，变量代换给出
\[
 \|D^\kappa u_\varepsilon\|_{L^p(B)}
       =\varepsilon^{k-|\kappa|}
                    \|D^\kappa\varphi\|_{L^p(\mathbb R^d)}.
\]
因此这列函数的 \(W^{k,p}(B)\) 范数一致有界；最高阶项保持不变，低阶项趋于零。
另一方面，对 \(q<\infty\)，
\[
 \|u_\varepsilon\|_{L^q(B)}
       =\varepsilon^{k-d/p+d/q}\|\varphi\|_q,
\]
而 \(q=\infty\) 时同式按 \(d/q=0\) 理解。若 \(kp<d\) 且
\(q>dp/(d-kp)\)，这个幂次严格为负，左端趋于无穷，与统一估计矛盾。

为检验第二问，选取 \(\varphi\)、一个 \(|\kappa|=j\) 的多重指标，以及
\(a,b\in B\)，使 \(a\ne b\) 且
\(D^\kappa\varphi(a)\ne D^\kappa\varphi(b)\)。例如取一个相应导数不恒定的光滑紧支集函数即可。对同一集中族，在两个点
\(\varepsilon a,\varepsilon b\) 处计算得到
\[
 \frac{|D^\kappa u_\varepsilon(\varepsilon a)
             -D^\kappa u_\varepsilon(\varepsilon b)|}
      {|\varepsilon a-\varepsilon b|^\beta}
 =\varepsilon^{k-d/p-j-\beta}
       \frac{|D^\kappa\varphi(a)-D^\kappa\varphi(b)|}
            {|a-b|^\beta}.
\]
最后一个分式是严格为正的固定数。因此若
\(j+\beta>k-d/p\)，左侧趋于无穷，而右侧所允许的 Sobolev 范数仍然有界，矛盾。

在临界关系 \(kp=d\) 下，第一问的函数族变成
\(u_\varepsilon(x)=\varphi(x/\varepsilon)\)，其上确界范数不随 \(\varepsilon\) 改变。这个计算既没有制造上确界发散，也没有证明统一 \(L^\infty\) 嵌入成立。临界情形需要对数剖面等不同于单纯重缩放的构造。
\end{solution}

\begin{exercise}\label{b1:ex:24-morrey-integrability-balance}
设 \(d<p<\infty\)、\(1\le r<\infty\)，令
\[
 \alpha=1-\frac dp,\quad
 a=\frac d{d+\alpha r}.
\]
证明对于每个
\(u\in W^{1,p}(\mathbb R^d)\cap L^r(\mathbb R^d)\)，其连续代表满足
\[
 \|u\|_\infty
    \le C_{d,p,r}\|\nabla u\|_p^a
                    \cdot\|u\|_r^{\,1-a}.
\]
提示：比较一个点的函数值与半径 \(R\) 的球平均，再依两个范数选择 \(R\)。
还应处理梯度范数为零的情形，并说明为何任何具有相同两因子形式、对全空间全部函数成立的估计，都必须采用这里的指数 \(a\)。
\end{exercise}
% 来源：自拟；将全空间 Morrey 振荡控制与额外 Lr 可积性平衡，不使用紧性或插值空间理论。
\begin{solution}
记 \(G=\|\nabla u\|_p\)、\(A=\|u\|_r\)。由定理%
\ref{b1:thm:morrey-whole-space}，存在唯一连续代表，使
\[
 |u(x)-u(y)|\le C_{d,p}G|x-y|^\alpha
                  \quad(x,y\in\mathbb R^d).
\]
若 \(G=0\)，这个代表为全空间常函数；有限 \(r\) 的可积性迫使它等于零，所求估计成立。
若 \(A=0\)，函数同样为零。因此以下只考虑 \(G,A>0\)。

固定 \(x\in\mathbb R^d\) 及任意 \(R>0\)，对 \(y\in B(x,R)\) 平均上述振荡界，再使用 Hölder 不等式，得到
\[
 \begin{aligned}
 |u(x)|
 &\le\frac1{|B(x,R)|}\int_{B(x,R)}|u(y)|\,\mathrm dy
                                  +C_{d,p}G R^\alpha\\
 &\le\omega_d^{-1/r}R^{-d/r}A+C_{d,p}G R^\alpha.
 \end{aligned}
\]
在 \(r=1\) 时，第二行只是把局部积分放大为全空间积分。选择
\[
 R=(A/G)^{1/(\alpha+d/r)}
\]
后，两个项都由一个只依赖 \(d,p,r\) 的常数乘
\[
 G^{\,d/(d+\alpha r)}A^{\,\alpha r/(d+\alpha r)}
\]
控制。此界与 \(x\) 无关，取上确界即得结论。

最后，假设同样的估计能以某个固定指数 \(b\) 代替 \(a\)。取非零
\(\varphi\in C_c^\infty(\mathbb R^d)\)，对全部 \(\lambda>0\) 测试
\(\varphi_\lambda(x)=\varphi(\lambda x)\)。其上确界不变，而
\[
 \|\nabla\varphi_\lambda\|_p
        =\lambda^\alpha\|\nabla\varphi\|_p,\quad
 \|\varphi_\lambda\|_r
        =\lambda^{-d/r}\|\varphi\|_r.
\]
估计右端中的尺度因子为 \(\lambda^{b\alpha-(1-b)d/r}\)。若幂次非零，令
\(\lambda\) 趋于零或无穷即可使右端趋于零，违反左端严格为正这一事实。所以
\(b\alpha=(1-b)d/r\)，解得 \(b=a\)。这一必要性来自尺度平衡，并未使用任何极值函数或紧性结论。
\end{solution}
